\documentclass[../main.tex]{subfiles}
\graphicspath{{\subfix{../}}}
\begin{document}

\textit{Csillapított és kényszerrezgések, csatolt rezgések, lineáris lánc. Kepler-probléma, bolygómozgás. Kvantummechanikai problémák: potenciálvölgy, oszcillátor, rotátor. Hidrogénatom. Keltő és eltüntető operátorok.}

\subsection{Csillapított és kényszerrezgések, csatolt rezgések, lineáris lánc}
\subsubsection{Csillapított rezgés}
Csillapított rezgésről akkor beszélünk, ha egy rezgő rendszerben van valamilyen energia veszteség, amely csökkenti a rezgés amplitúdóját idővel. A csillapítás lehet lineáris vagy nemlineáris, és különböző fizikai mechanizmusok révén jöhet létre,
például súrlódás, légellenállás vagy belső súrlódás. Vegyünk példának egy olyan rendszert amikor a harmonikus oszcillátor egy közegben mozog, amely a sebességgel arányos módon csillapítja a mozgást:
\begin{figure}[H]
    \centering
    \includegraphics[width=0.4\textwidth]{imgs/3-tetel/csillapitott_rezges.jpg}
    \caption{Harmonikus oszcillátor közegben}
    \label{fig:csillapitott_rezges}
\end{figure}
Ekkor a mozgásegyenlet a következő lesz:
\begin{equation}
    m\ddot{x} = - kx - b\dot{x}
\end{equation}
Ahol $b$ a csillapítási együttható, amely a közeg tulajdonságaitól függ. Itt látható, hogy a standard harmonikus oszcillátorhoz képest van egy plusz tag, amely a sebességgel arányos és negatív előjelű, tehát csökkenti a mozgást. Ez egy másodrendű lineáris differenciálegyenlet,
amelyet megoldhatunk a klasszikus Ansatz segítségével. Először is osszunk le $m$-el és paraméterezzük át a mozgásegyenletet:
\begin{equation}
    \ddot{x} + 2\beta \dot{x} + \omega_0^2 x = 0
\end{equation}
\begin{equation}
    \text{Ahol } \beta = \frac{b}{2m} \quad \text{és} \quad \omega_0 = \sqrt{\frac{k}{m}}
\end{equation}
Mivel egy periodikus függvényt feltételezünk ezért próbálkozzunk meg egy $\sin$ fügvénnyel:
\begin{equation}
    x(t) = A(t) sin(\omega t + \varphi)
\end{equation}
Nézzük meg, hogy ez a függvény kielégíti-e a mozgásegyenletet. Ehhez szükségünk lesz az első és második deriváltjára:
\begin{equation}
    \dot{x}(t) = \dot{A}(t) sin(\omega t + \varphi) + A(t) \omega cos(\omega t + \varphi)
\end{equation}
\begin{equation}
    \ddot{x}(t) = \ddot{A}(t) sin(\omega t + \varphi) + 2\dot{A}(t) \omega cos(\omega t + \varphi) - A(t) \omega^2 sin(\omega t + \varphi)
\end{equation}
Ezt behelyettesítve a mozgásegyenletbe:
\begin{equation}
    \begin{aligned}
        & \ddot{A}(t) sin(\omega t + \varphi) + 2\dot{A}(t) \omega cos(\omega t + \varphi) - A(t) \omega^2 sin(\omega t + \varphi) + \\
        & 2\beta \left( \dot{A}(t) sin(\omega t + \varphi) + A(t) \omega cos(\omega t + \varphi) \right) + \\
        & \omega_0^2 A(t) sin(\omega t + \varphi) = 0
    \end{aligned}
\end{equation}
Most csoportosítsuk a $sin$ és $cos$ tagokat külön:
\begin{equation}
    \left( \ddot{A}(t) - A(t) \omega^2 + 2\beta \dot{A}(t) + \omega_0^2 A(t) \right) sin(\omega t + \varphi) + \left( 2\dot{A}(t) \omega + 2\beta A(t) \omega \right) cos(\omega t + \varphi) = 0
\end{equation}
Mivel a $sin$ és $cos$ függvények lineárisan függetlenek, ezért mindkét zárójelnek nullának kell lennie:
\begin{equation}
    \ddot{A}(t) - A(t) \omega^2 + 2\beta \dot{A}(t) + \omega_0^2 A(t) = 0
\end{equation}
\begin{equation}
    2\dot{A}(t) \omega + 2\beta A(t) \omega = 0
\end{equation}
Nézzük meg először az elsőrendű egyenletet:
\begin{equation}
    \dot{A}(t) + \beta A(t) = 0
\end{equation}
Ennek a megoldása triviálisan:
\begin{equation}
    A(t) = A_0 e^{-\beta t}
\end{equation}
Ezt behelyettesítve a másodrendű egyenletbe:
\begin{equation}
    \beta^2 A_0 e^{-\beta t} - A_0 e^{-\beta t} \omega^2 - 2\beta^2 A_0 e^{-\beta t} + \omega_0^2 A_0 e^{-\beta t} = 0
\end{equation}
\begin{equation}
    -\beta^2 - \omega^2 + \omega_0^2 = 0
\end{equation}
\begin{equation}
    \omega^2 = \omega_0^2 - \beta^2
\end{equation}
Tehát a csillapított rezgés megoldása a következő:
\begin{equation}
    x(t) = A_0 e^{-\beta t} sin(\omega t + \varphi)
\end{equation}
Ahol $\omega = \sqrt{\omega_0^2 - \beta^2}$. Ez azt jelenti, hogy a mozgás nagyban függ a kezdeti feltételektől, amik alapján három esetet különböztetünk meg:
\begin{itemize}
    \item Alulcsillapított rezgés ($\beta < \omega_0$): Ilyenkor a rendszer rezeg, de az amplitúdó exponenciálisan csökken idővel. A rezgés frekvenciája kisebb, mint a természetes frekvencia.
    \item Kritikus csillapítás ($\beta = \omega_0$): Ilyenkor a rendszer nem rezeg, hanem a lehető legy gyorsabban tér vissza az egyensúlyi helyzetbe anélkül, hogy túllépné azt.
    \item Túlcsillapított rezgés ($\beta > \omega_0$): Ilyenkor a rendszer szintén nem rezeg, de lassabban tér vissza az egyensúlyi helyzetbe, mint a kritikus csillapítás esetén.
\end{itemize}   
Tehát minden esetben exponenciálisan csökken az amplitúdó, de a mozgás jellege és a visszatérés sebessége az egyensúlyi helyzetbe a csillapítás mértékétől függ.
\begin{figure}[H]
\centering
\begin{subfigure}{.5\textwidth}
  \centering
  \includegraphics[width=1\linewidth]{imgs/3-tetel/csillapitas_merteke.jpg}
  \caption{Alulcsillapított rezgés alakja}
  \label{fig:alulcsillapitott}
\end{subfigure}%
\begin{subfigure}{.5\textwidth}
  \centering
  \includegraphics[width=1\linewidth]{imgs/3-tetel/csillapitott_rezgesek_tipusai.png}
  \caption{Csillapított rezgések típusai}
  \label{fig:kritikus_csillapitott}
\end{subfigure}
\caption{Csillapított rezgések}
\label{fig:csillapitott_rezgesek}
\end{figure}

\subsubsection{kényszerrezgések}
A kényszerrezgés bizonyos szempontból a Csillapított rezgés ellentéte, hiszen itt egy külső erő hat a rendszerre, amely folyamatosan energiát visz be a rendszerbe, így fenntartva a rezgést.
A kényszerrezgés akkor fordul elő, amikor egy rezgő rendszerre egy időben változó külső erő hat, amelynek frekvenciája közel van a rendszer sajátfrekvenciájához.
Ez a jelenség rezonanciaként ismert, és akkor következik be, amikor a külső erő frekvenciája megegyezik a rendszer természetes frekvenciájával, vagy nagyon közel van hozzá. Ilyenkor a rendszer válasza jelentősen megnő, és az amplitúdója is nagyobb lesz.
Vegyünk egy példát egy harmonikus oszcillátorra, amelyre egy külső erő hat:
\begin{figure}[H]
    \centering
    \includegraphics[width=0.8\textwidth]{imgs/3-tetel/kenyszerrezges.png}
    \caption{Harmonikus oszcillátor külső erő hatására}
    \label{fig:kenyszerrezges}
\end{figure}
írjuk fel a mozgásegyenletet:
\begin{equation}
    m\ddot{x} + b\dot{x} + kx = F_0 sin(\omega t)
\end{equation}
Ahol $F_0$ a külső erő amplitúdója, $\omega$ pedig a külső erő frekvenciája. Ebben az egyenletben jelen van a csillapítás is, amelyet a $b\dot{x}$ tag képvisel. Ez egy inhomogén másodrendű lineáris differenciálegyenlet,
amelynek a megoldása a homogén egyenlet megoldásából és az inhomogén megoldásból áll. Rendezzük át megint az egyenletet és paraméterezzük át:
\begin{equation}
    \ddot{x} + 2\beta \dot{x} + \omega_0^2 x = \frac{F_0}{m} sin(\omega t) 
\end{equation}
A homogén megoldás megegyezik a csillapított rezgés megoldásával, amelyet már korábban megtaláltunk:
\begin{equation}
    x_h(t) = A_0 e^{-\beta t} sin(\omega_d t + \varphi)
\end{equation}
Ahol $\omega_d = \sqrt{\omega_0^2 - \beta^2}$. Most keressük meg az inhomogén megoldást. Mivel a jobb oldalon egy szinuszos kényszerítő erő van, ezért próbálkozzunk meg egy szinuszos megoldással:
\begin{equation}
    x_p(t) = A sin(\omega t + \phi)
\end{equation}
Ezt behelyettesítve az egyenletbe:
\begin{equation}
    -A \omega^2 sin(\omega t + \phi) + 2\beta A \omega cos(\omega t + \phi) + \omega_0^2 A sin(\omega t + \phi) = \frac{F_0}{m} sin(\omega t)
\end{equation}
Csoportosítsuk a $sin$ és $cos$ tagokat külön:
\begin{equation}
    \left( -A \omega^2 + \omega_0^2 A \right) sin(\omega t + \phi) + 2\beta A \omega cos(\omega t + \phi) = \frac{F_0}{m} sin(\omega t)
\end{equation}
Most használjuk a szögösszeg képletet, hogy kifejezzük a $sin(\omega t + \phi)$ és $cos(\omega t + \phi)$-t:
\begin{equation}
    \left( -A \omega^2 + \omega_0^2 A \right) \left( sin(\omega t) cos(\phi) + cos(\omega t) sin(\phi) \right) + 2\beta A \omega \left( cos(\omega t) cos(\phi) - sin(\omega t) sin(\phi) \right) = \frac{F_0}{m} sin(\omega t)
\end{equation}
Csoportosítsuk a $sin(\omega t)$ és $cos(\omega t)$ tagokat külön:
\begin{equation}
    \begin{aligned}
        & \left( \left( -A \omega^2 + \omega_0^2 A \right) cos(\phi) - 2\beta A \omega sin(\phi) \right) sin(\omega t) + \\
        & \left( \left( -A \omega^2 + \omega_0^2 A \right) sin(\phi) + 2\beta A \omega cos(\phi) \right) cos(\omega t) = \frac{F_0}{m} sin(\omega t)
    \end{aligned}
\end{equation}
Mivel a $sin(\omega t)$ és $cos(\omega t)$ függvények lineárisan függetlenek, ezért a cos-os tagnak el kell tűnnie, és a sin-es tagnak meg kell egyeznie a jobb oldallal:
\begin{equation}
    \left( -A \omega^2 + \omega_0^2 A \right) cos(\phi) - 2\beta A \omega sin(\phi) = \frac{F_0}{m}
\end{equation}
\begin{equation}
    \left( -A \omega^2 + \omega_0^2 A \right) sin(\phi) + 2\beta A \omega cos(\phi) = 0
\end{equation}
Itt a két egyenletet négzetre emelve és összeadva kifejezhetjük az $A$ amplitúdót:
\begin{equation}
    A^2\left[\left(\left( -\omega^2 + \omega_0^2 \right) cos(\phi) - 2\beta  \omega sin(\phi)\right)^2 + \left( \left( - \omega^2 + \omega_0^2 \right) sin(\phi) + 2\beta \omega cos(\phi) \right)^2\right] = \left(\frac{F_0}{m}\right)^2
\end{equation}
\begin{equation}
    \begin{aligned}
        & A^2 \left(-\omega^2 + \omega^2_0\right)^2 cos^2(\phi) + \left(2\beta \omega\right)^2 sin^2(\phi) - 4\beta \omega \left(-\omega^2 + \omega^2_0\right) sin(\phi) cos(\phi) + \\
        & A^2 \left(-\omega^2 + \omega^2_0\right)^2 sin^2(\phi) + \left(2\beta \omega\right)^2 cos^2(\phi) + 4\beta \omega \left(-\omega^2 + \omega^2_0\right) sin(\phi) cos(\phi) = \left(\frac{F_0}{m}\right)^2
    \end{aligned}
\end{equation}
\begin{equation}
    A^2 \left[ \left(-\omega^2 + \omega^2_0\right)^2 \left( cos^2(\phi) + sin^2(\phi) \right) + \left(2\beta \omega\right)^2 \left( sin^2(\phi) + cos^2(\phi) \right) \right] = \left(\frac{F_0}{m}\right)^2
\end{equation}
\begin{equation}
    A^2 \left[ \left(-\omega^2 + \omega^2_0\right)^2 + \left(2\beta \omega\right)^2 \right] = \left(\frac{F_0}{m}\right)^2
\end{equation}
\begin{equation}
    A = \frac{\frac{F_0}{m}}{\sqrt{(\omega_0^2 - \omega^2)^2 + (2\beta \omega)^2}}
\end{equation}
Majd a második egyenletből kifejezhetjük a $\phi$ fázist:
\begin{equation}
    A\left[\left( - \omega^2 + \omega_0^2  \right) sin(\phi) + 2\beta  \omega cos(\phi)\right] = 0
\end{equation}
\begin{equation}
    \left( - \omega^2 + \omega_0^2  \right) sin(\phi) + 2\beta  \omega cos(\phi) = 0
\end{equation}
\begin{equation}
    \left( - \omega^2 + \omega_0^2  \right) sin(\phi) = - 2\beta  \omega cos(\phi)
\end{equation}
\begin{equation}
    tan(\phi) = \frac{- 2\beta  \omega}{- \omega^2 + \omega_0^2}
\end{equation}
\begin{equation}
    tan(\phi) = \frac{2\beta  \omega}{\omega^2 - \omega_0^2}
\end{equation}
Tehát a kényszerrezgés teljes megoldása a homogén és inhomogén megoldás összege:
\begin{equation}
    x(t) = A_0 e^{-\beta t} sin(\omega_d t + \varphi) + \frac{\frac{F_0}{m}}{\sqrt{(\omega_0^2 - \omega^2)^2 + (2\beta \omega)^2}} sin(\omega t + \phi)
\end{equation}
Ahol az első tag a csillapított rezgés, a második tag pedig a kényszerrezgés. Idővel a csillapított rezgés eltűnik, és a rendszer csak a kényszerrezgést mutatja. Az amplitúdó a kényszerrezgésnél a következőképpen viselkedik:
\begin{itemize}
    \item Ha a kényszerítő erő frekvenciája nagyon eltér a rendszer sajátfrekvenciájától ($\omega \ll \omega_0$ vagy $\omega \gg \omega_0$), akkor az amplitúdó kicsi lesz, mivel a rendszer nem rezonál a kényszerítő erővel.
    \item Ha a kényszerítő erő frekvenciája közel van a rendszer sajátfrekvenciájához ($\omega \approx \omega_0$), akkor az amplitúdó jelentősen megnő, mivel a rendszer rezonál a kényszerítő erővel. Ez a rezonancia jelensége.
    \item Ha a csillapítás kicsi ($\beta \to 0$), akkor az amplitúdó nagyon nagy lehet rezonancia esetén, mivel nincs elég energia veszteség a rendszerben.
    \item Ha a csillapítás nagy ($\beta$ nagy), akkor az amplitúdó kisebb lesz rezonancia esetén is, mivel a csillapítás elnyeli az energiát, és megakadályozza, hogy az amplitúdó túl nagy legyen.
\end{itemize}
\begin{figure}[H]
    \centering
    \includegraphics[width=0.6\textwidth]{imgs/3-tetel/resonance.png}
    \caption{Kényszerrezgés amplitúdója a kényszerítő erő frekvenciájának függvényében különböző csillapítási értékek mellett}
    \label{fig:rezges_amplitudo}
\end{figure}

\subsubsection{Csatolt rezgések}
A csatolt rezgések olyan rendszerek, amelyekben két vagy több rezgő elem kölcsönhatásban áll egymással, így a rezgésük nem független, hanem összefügg. A csatolt rezgések gyakran előfordulnak fizikai rendszerekben, például mechanikai rendszerekben,
elektromos rendszerekben vagy akár biológiai rendszerekben is. Vegyünk egy egyszerű példát két csatolt harmonikus oszcillátorra, amelyek egy rugóval vannak összekötve:
\begin{figure}[H]
    \centering
    \includegraphics[width=0.6\textwidth]{imgs/3-tetel/csatolt_rezges.png}
    \caption{Két csatolt harmonikus oszcillátor}
    \label{fig:csatolt_rezges}
\end{figure}
A mozgásegyenletek a következők lesznek:
\begin{equation}
    m\ddot{x}_1 = -k x_1 - k_c (x_1 - x_2)
\end{equation}
\begin{equation}
    m\ddot{x}_2 = -k x_2 - k_c (x_2 - x_1)
\end{equation}
Ahol $k$ a rugóállandó, $k_c$ a csatolási rugóállandó, $m$ a tömeg, $x_1$ és $x_2$ pedig az első és második oszcillátor helyzete. Ezeket az egyenleteket rendezzük át:
\begin{equation}
    \ddot{x}_1 + \left(\frac{k + k_c}{m}\right) x_1 - \frac{k_c}{m} x_2 = 0
\end{equation}
\begin{equation}
    \ddot{x}_2 + \left(\frac{k + k_c}{m}\right) x_2 - \frac{k_c}{m} x_1 = 0
\end{equation}
Nevezzük át ismét a paramétereket:
\begin{equation}
    \omega_0^2 = \frac{k + k_c}{m} \quad \quad \omega_c^2 = \frac{k_c}{m}
\end{equation}
Így az egyenletek a következőképpen néznek ki:
\begin{equation}
    \ddot{x}_1 + \omega_0^2 x_1 - \omega_c^2 x_2 = 0
\end{equation}
\begin{equation}
    \ddot{x}_2 + \omega_0^2 x_2 - \omega_c^2 x_1 = 0
\end{equation}
Adjuk össze és vonjuk ki az egyenleteket:
\begin{equation}
    \ddot{x}_1 + \ddot{x}_2 + \omega_0^2 (x_1 + x_2) - \omega_c^2 (x_2 + x_1) = 0
\end{equation}
\begin{equation}
    \ddot{x}_1 - \ddot{x}_2 + \omega_0^2 (x_1 - x_2) - \omega_c^2 (x_2 - x_1) = 0
\end{equation}
Nevezzük el az új változókat:
\begin{equation}
    X = x_1 + x_2 \quad \quad Y = x_1 - x_2
\end{equation}
Így az egyenletek a következőképpen néznek ki:
\begin{equation}
    \ddot{X} + (\omega_0^2 - \omega_c^2) X = 0
\end{equation}
\begin{equation}
    \ddot{Y} + (\omega_0^2 + \omega_c^2) Y = 0
\end{equation}
Ezek már független egyenletek, amelyeket könnyen megoldhatunk:
\begin{equation}
    \ddot{X} = - (\omega_0^2 - \omega_c^2) X
\end{equation}
\begin{equation}
    \ddot{Y} = - (\omega_0^2 + \omega_c^2) Y
\end{equation}
Innen triviálisan megkapjuk a megoldásokat:
\begin{equation}
    X(t) = A_1 sin(\omega_1 t + \varphi_1) \quad \quad \omega_1 = \sqrt{\omega_0^2 - \omega_c^2}
\end{equation}
\begin{equation}
    Y(t) = A_2 sin(\omega_2 t + \varphi_2) \quad \quad \omega_2 = \sqrt{\omega_0^2 + \omega_c^2}
\end{equation}
Visszaalakítva az eredeti változókra:
\begin{equation}
    x_1(t) = \frac{X(t) + Y(t)}{2} = \frac{A_1 sin(\omega_1 t + \varphi_1) + A_2 sin(\omega_2 t + \varphi_2)}{2}
\end{equation}
\begin{equation}
    x_2(t) = \frac{X(t) - Y(t)}{2} = \frac{A_1 sin(\omega_1 t + \varphi_1) - A_2 sin(\omega_2 t + \varphi_2)}{2}
\end{equation}
Ahol $\omega_1$ és $\omega_2$ a két csatolt rezgés sajátfrekvenciái, amelyek a csatolás mértékétől függnek. Az $A_1$, $A_2$, $\varphi_1$ és $\varphi_2$ a kezdeti feltételektől függő állandók.

Nézzünk meg egy speciális esetet, amikor $A_1 = A_2 = \frac{A}{2}$ és $\varphi_1 = \varphi_2 = \frac{\pi}{2}$ (azaz mindkét oszcillátor ugyanazzal az amplitúdóval és fázissal indul):
\begin{equation}
    x_1(t) = \frac{A}{2} \left( sin(\omega_1 t + \frac{\pi}{2}) + sin(\omega_2 t + \frac{\pi}{2}) \right) = \frac{A}{2} \left( cos(\omega_1 t) + cos(\omega_2 t) \right)
\end{equation}
\begin{equation}
    x_2(t) = \frac{A}{2} \left( sin(\omega_1 t + \frac{\pi}{2}) - sin(\omega_2 t + \frac{\pi}{2}) \right) = \frac{A}{2} \left( cos(\omega_1 t) - cos(\omega_2 t) \right)
\end{equation}
A $cos(x) + cos(y) = 2 cos\left(\frac{x+y}{2}\right) cos\left(\frac{x-y}{2}\right)$ és $cos(x) - cos(y) = -2 sin\left(\frac{x+y}{2}\right) sin\left(\frac{x-y}{2}\right)$ azonosságot felhasználva:
\begin{equation}
    x_1(t) = A cos\left(\frac{(\omega_1 + \omega_2)t}{2}\right) cos\left(\frac{(\omega_1 - \omega_2)t}{2}\right)
\end{equation}
\begin{equation}
    x_2(t) = A sin\left(\frac{(\omega_1 + \omega_2)t}{2}\right) sin\left(\frac{(\omega_1 - \omega_2)t}{2}\right)
\end{equation}
Ha emellé feltesszük, hogy a csatolás kicsi, azaz $\omega_1 \approx \omega_2$, akkor a sajátfrekvenciák közel lesznek egymáshoz:
\begin{equation}
    \omega = \frac{\omega_1 + \omega_2}{2} \approx \omega_1 \approx \omega_2 \quad \quad \epsilon = \frac{\omega_1 - \omega_2}{2} \ll \omega
\end{equation}
Így a megoldások a következőképpen egyszerűsödnek:
\begin{equation}
    x_1(t) \approx A cos(\omega t) cos\left(\frac{(\omega_1 - \omega_2)t}{2}\right)
\end{equation}
\begin{equation}
    x_2(t) \approx A sin(\omega t) sin\left(\frac{(\omega_1 - \omega_2)t}{2}\right)
\end{equation}
Ez azt jelenti, hogy mindkét oszcillátor rezeg a $\omega_0$ frekvenciával, de az amplitúdójuk időben változik, és egy lassú modulációt mutat a $\frac{\omega_1 - \omega_2}{2}$ frekvenciával. Ez a jelenség a lebegés a csatolt rezgések egyik jellegzetessége,
és gyakran előfordul különböző fizikai rendszerekben, például molekulákban, kristályokban vagy akár mechanikai rendszerekben is.
\begin{figure}[H]
    \centering
    \includegraphics[width=0.6\textwidth]{imgs/3-tetel/lebeges.png}
    \caption{Lebegés jelensége csatolt rezgések esetén}
    \label{fig:csatolt_rezgesek_idobeli}
\end{figure}

\subsubsection{Lineáris lánc}
A lineáris lánc a csatolt rezgés kiterjesztése egy végtelen számú oszcillátorra, amelyek egy egyenes vonal mentén helyezkednek el, és csak a szomszédos oszcillátorokkal vannak csatolva. Ez a modell gyakran használatos a kristályrácsok vibrációinak leírására,
mivel a kristályok atomjai hasonlóan rendeződnek el egy rácsban, és csak a közvetlenül szomszédos atomokkal lépnek kölcsönhatásba. Vegyünk egy végtelen lineáris láncot, ahol minden oszcillátor tömege $m$, és a csatolási rugóállandó $k$. Emelett feltételezzünk
periodikus határfeltételeket, azaz az első és utolsó oszcillátor is csatolva van egymáshoz, így a határokon is tudjuk kényelmesen kezelni a problémát:
\begin{figure}[H]
    \centering
    \includegraphics[width=0.8\textwidth]{imgs/3-tetel/linearis_lanc.png}
    \caption{Végtelen lineáris lánc}
    \label{fig:linear_chain}
\end{figure}
A mozgásegyenlet az n-edik oszcillátor számára a következő lesz:
\begin{equation}
    m\ddot{x}_n = -k (x_n - x_{n-1}) - k (x_n - x_{n+1})
\end{equation}
Ahol $x_n$ az $n$-edik oszcillátor helyzete. Ezt rendezzük át:
\begin{equation}
    \ddot{x}_n = -\frac{k}{m} (2x_n - x_{n-1} - x_{n+1})
\end{equation}
Nevezzük át a paramétert:
\begin{equation}
    \omega_0^2 = \frac{k}{m}
\end{equation}
Így az egyenlet a következőképpen néz ki:
\begin{equation}
    \ddot{x}_n = -\omega_0^2 (2x_n - x_{n-1} - x_{n+1})
\end{equation}
Hogy az összes mozgásegyenletet egyszerre tudjuk kezelni, tegyük egy mátrixba az összes $x_n$-t:
\begin{equation}
    \mathbf{x} = \begin{pmatrix}
    \vdots \\
    x_{n-1} \\
    x_n \\
    x_{n+1} \\
    \vdots
    \end{pmatrix}
\end{equation}
Így a mozgásegyenletek mátrix formában a következőképpen néznek ki:
\begin{equation}
    \ddot{\mathbf{x}} = -\omega_0^2 \mathbf{\hat{D}x}
\end{equation}
Ahol $\mathbf{\hat{D}}$ a diszkrét második differenciál operátora, amely a következőképpen néz ki:
\begin{equation}
    \mathbf{\hat{D}} = \begin{pmatrix}
    - 2 & 1 & 0 & 0 & \cdots & 1 \\
    1 & -2 & 1 & 0 & \cdots & 0 \\
    0 & 1 & -2 & 1 & \cdots & 0 \\
    0 & 0 & 1 & -2 & \cdots & 0 \\
    \vdots & \vdots & \vdots & \vdots & \ddots & \vdots \\
    1 & 0 & 0 & 0 & \cdots & -2
    \end{pmatrix}
\end{equation}
Keressük a megoldást periodikus hullámok formájában:
\begin{equation}
    x_n(t) = A_n(t) e^{i\omega t}
\end{equation}
Továbbá tudjuk, hogy az amplitúdó is periodikus, a csatolt rezgés alapján:
\begin{equation}
    A_n(t) = A e^{i q n a}
\end{equation}
Ahol $q$ a hullámszám, $a$ pedig az oszcillátorok közötti távolság. Így a teljes megoldás:
\begin{equation}
    x_n(t) = A e^{i (q n a + \omega t)}
\end{equation}
Ezt behelyettesítve a mozgásegyenletbe:
\begin{equation}
    -\omega^2 A e^{i (q n a + \omega t)} = -\omega_0^2 \left( 2A e^{i (q n a + \omega t)} - A e^{i (q (n-1) a + \omega t)} - A e^{i (q (n+1) a + \omega t)} \right)
\end{equation}
Egyszerűsítve:
\begin{equation}
    -\omega^2 = -\omega_0^2 \left( 2 - e^{-i q a} - e^{i q a} \right)
\end{equation}
\begin{equation}
    \omega^2 = \omega_0^2 \left( 2 - 2 cos(q a) \right)
\end{equation}
felhasználva a $1 - cos(x) = 2 sin^2\left(\frac{x}{2}\right)$ azonosságot:
\begin{equation}
    \omega^2 = 4 \omega_0^2 sin^2\left(\frac{q a}{2}\right)
\end{equation}
\begin{equation}
    \omega = 2 \omega_0 \left| sin\left(\frac{q a}{2}\right) \right|
\end{equation}
Ezt nevezik diszperziós relációnak, amely leírja a hullámok frekvenciáját a hullámszám függvényében egy végtelen lineáris láncban. Ez a reláció fontos információkat nyújt a rendszer dinamikájáról,
például a hullámok terjedési sebességéről és a rendszer rezgési módjairól. A diszperziós reláció alapján látható, hogy a frekvencia $\omega$ a hullámszám $q$ függvényében szinuszosan változik, és a maximális frekvencia
$\omega_{max} = 2\omega_0$ akkor érhető el, amikor $q a = \pi$. Ez azt jelenti, hogy a hullámok nem terjedhetnek végtelen sebességgel a lineáris láncban, hanem van egy felső határa frekvenciára, amelyet a rendszer fizikai tulajdonságai határoznak meg.
Emellett a reciprokrács tulajdonságai miatt a hullámszám $q$ értékei csak egy bizonyos tartományban értelmezettek, amelyet Brillouin-zónának neveznek. Az első Brillouin-zóna a $-\frac{\pi}{a} \leq q \leq \frac{\pi}{a}$ intervallumot fedi le, és ebben a tartományban
minden lehetséges hullámszám érték megtalálható. Ez  a zóna fontos szerepet játszik a kristályok és más periodikus struktúrák fizikai tulajdonságainak megértésében, mivel a hullámok terjedése és kölcsönhatása a kristályrács szerkezetével
határozza meg a rendszer viselkedését.
\begin{figure}[H]
    \centering
    \includegraphics[width=0.4\textwidth]{imgs/3-tetel/egyatomos_linearis_lanc.png}
    \caption{Példa egyatomos lineáris lánc diszperziós relációjára}
    \label{fig:linearis_lanc_diszperzio}
\end{figure}

\subsubsection*{Példa: kétatomos lineáris lánc}
A kétatomos lineáris lánc egy olyan modell, amelyben két különböző tömegű részecske (atom) váltakozva helyezkedik el egy egyenes vonal mentén, és csak a szomszédos atomokkal vannak csatolva. Ez a modell gyakran használatos a kristályrácsok vibrációinak
leírására, különösen olyan anyagok esetében, ahol két különböző típusú atom található, például sókristályokban (pl. NaCl). Vegyünk egy végtelen kétatomos lineáris láncot, ahol az egyik atom tömege $m_1$, a másik atom tömege $m_2$, és a csatolási rugóállandó $k$. Feltételezzünk
periodikus határfeltételeket, azaz az első és utolsó atom is csatolva van egymáshoz:
\begin{figure}[H]
    \centering
    \includegraphics[width=0.8\textwidth]{imgs/3-tetel/ketatomos_linearis_lanc.png}
    \caption{Végtelen kétatomos lineáris lánc}
    \label{fig:ketoatomos_linear_chain}
\end{figure}
A mozgásegyenletek az n-edik pár atom számára a következő lesz:
\begin{equation}
    m_1 \ddot{x}_n = -k (x_n - y_n) - k (x_n - y_{n-1})
\end{equation}
\begin{equation}
    m_2 \ddot{y}_n = -k (y_n - x_n) - k (y_n - x_{n+1})
\end{equation}
Ahol $x_n$ az n-edik $m_1$ tömegű atom helyzete, $y_n$ pedig az n-edik $m_2$ tömegű atom helyzete. Ezt rendezzük át:
\begin{equation}
    \ddot{x}_n = -\frac{k}{m_1} (2x_n - y_n - y_{n-1})
\end{equation}
\begin{equation}
    \ddot{y}_n = -\frac{k}{m_2} (2y_n - x_n - x_{n+1})
\end{equation}
Nevezzük át a paramétereket:
\begin{equation}
    \omega_1^2 = \frac{k}{m_1} \quad \quad \omega_2^2 = \frac{k}{m_2}
\end{equation}
Így az egyenletek a következőképpen néznek ki:
\begin{equation}
    \ddot{x}_n = -\omega_1^2 (2x_n - y_n - y_{n-1})
\end{equation}
\begin{equation}
    \ddot{y}_n = -\omega_2^2 (2y_n - x_n - x_{n+1})
\end{equation}
Keressük a megoldást periodikus hullámok formájában:
\begin{equation}
    x_n(t) = A e^{i (q n a + \omega t)}
\end{equation}
\begin{equation}
    y_n(t) = B e^{i (q n a + \omega t)}
\end{equation}
Ahol $q$ a hullámszám, $a$ pedig az atomok közötti távolság. Ezt behelyettesítve a mozgásegyenletekbe:
\begin{equation}
    -\omega^2 A e^{i (q n a + \omega t)} = -\omega_1^2 \left( 2A e^{i (q n a + \omega t)} - B e^{i (q n a + \omega t)} - B e^{i (q (n-1) a + \omega t)} \right)
\end{equation}
\begin{equation}
    -\omega^2 B e^{i (q n a + \omega t)} = -\omega_2^2 \left( 2B e^{i (q n a + \omega t)} - A e^{i (q n a + \omega t)} - A e^{i (q (n+1) a + \omega t)} \right)
\end{equation}
Egyszerűsítve:
\begin{equation}
    -\omega^2 A = -\omega_1^2 \left( 2A - B - B e^{-i q a} \right)
\end{equation}
\begin{equation}
    -\omega^2 B = -\omega_2^2 \left( 2B - A - A e^{i q a} \right)
\end{equation}
Ezeket az egyenleteket mátrix formában is felírhatjuk:
\begin{equation}
    \begin{pmatrix}
    \omega^2 - 2\omega_1^2 & \omega_1^2 (1 + e^{-i q a}) \\
    \omega_2^2 (1 + e^{i q a}) & \omega^2 - 2\omega_2^2
    \end{pmatrix}
    \begin{pmatrix}
    A \\
    B
    \end{pmatrix}
    = 0
\end{equation}
Ahol a mátrix együtthatóinak determinánsának nullának kell lennie, hogy ne triviális megoldást kapjunk:
\begin{equation}
    \left| \begin{matrix}
    \omega^2 - 2\omega_1^2 & \omega_1^2 (1 + e^{-i q a}) \\
    \omega_2^2 (1 + e^{i q a}) & \omega^2 - 2\omega_2^2
    \end{matrix} \right| = 0
\end{equation}
Kiszámolva a determinánst:
\begin{equation}
    (\omega^2 - 2\omega_1^2)(\omega^2 - 2\omega_2^2) - \omega_1^2 \omega_2^2 (1 + e^{-i q a})(1 + e^{i q a}) = 0
\end{equation}
\begin{equation}
    (\omega^2 - 2\omega_1^2)(\omega^2 - 2\omega_2^2) - \omega_1^2 \omega_2^2 (2 + 2 cos(q a)) = 0
\end{equation}
\begin{equation}
    \omega^4 - 2(\omega_1^2 + \omega_2^2) \omega^2 + 4\omega_1^2 \omega_2^2 - 2\omega_1^2 \omega_2^2 (1 + cos(q a)) = 0
\end{equation}
\begin{equation}
    \omega^4 - 2(\omega_1^2 + \omega_2^2) \omega^2 + 2\omega_1^2 \omega_2^2 (1 - cos(q a)) = 0
\end{equation}
Ez egy másodfokú egyenlet $\omega^2$-re, amelynek megoldása a következő:
\begin{equation}
    \omega^2 = (\omega_1^2 + \omega_2^2) \pm \sqrt{(\omega_1^2 + \omega_2^2)^2 - 2\omega_1^2 \omega_2^2 (1 - cos(q a))}
\end{equation}
\begin{equation}
    \omega^2 = (\omega_1^2 + \omega_2^2) \pm \sqrt{(\omega_1^2 - \omega_2^2)^2 + 4\omega_1^2 \omega_2^2 sin^2\left(\frac{q a}{2}\right)}
\end{equation}
Ez a diszperziós reláció a kétatomos lineáris lánc számára. Két különböző ágat ad meg a frekvenciára, amelyeket akusztikus és optikai ágaknak neveznek:
\begin{itemize}
    \item Az akusztikus ág ($-$ jel) alacsony frekvenciákat ír le, ahol a két atom mozgása közel azonos fázisban történik. Ez az ág a hanghullámok terjedését modellezi a kristályban.
    \item Az optikai ág ($+$ jel) magas frekvenciákat ír le, ahol a két atom mozgása ellenfázisban történik. Ez az ág a fényhullámok kölcsönhatását modellezi a kristályban.
\end{itemize}
\begin{figure}[H]
    \centering
    \includegraphics[width=0.4\textwidth]{imgs/3-tetel/ketatomos_linearis_lanc_diszperzio.png}
    \caption{Példa kétatomos lineáris lánc diszperziós relációjára}
    \label{fig:ketatomos_linearis_lanc_diszperzio}
\end{figure}

\subsection{Kepler-Probléma, bolygómozgás}
A Kepler-probléma a klasszikus mechanikában egy olyan probléma, amely a két test közötti gravitációs kölcsönhatást vizsgálja, például egy bolygó és egy csillag között. A probléma célja meghatározni a bolygó pályáját és mozgását a csillag körül,
figyelembe véve a gravitációs erőt. A Kepler-probléma megoldása a Kepler-törvényeken alapul, amelyek Johannes Kepler német csillagász által a 17. században megfogalmazott három törvény:
\begin{itemize}
    \item \textbf{Első törvény (pályaalak törvénye)}: A bolygók ellipszis alakú pályán keringenek a Nap körül, amelynek egyik gyújtópontjában a Nap helyezkedik el.
    \item \textbf{Második törvény (területi sebesség törvénye)}: A bolygó és a Nap közötti képzeletbeli vonal egyenlő területeket súrol egyenlő idő alatt. Ez azt jelenti, hogy a bolygó gyorsabban mozog, amikor közelebb van a Naphoz, és lassabban, amikor távolabb van tőle.
    \item \textbf{Harmadik törvény (harmonikus törvény)}: A bolygók keringési idejének négyzete arányos a pálya fél nagytengelyének köbével. Matematikailag ez a következőképpen írható fel: $T^2 \propto a^3$, ahol $T$ a keringési idő, és $a$ a pálya fél nagytengelye.
\end{itemize}
A Kepler-probléma megoldása a Newton-féle gravitációs törvényen alapul, amely szerint a két test közötti gravitációs erő arányos a tömegeik szorzatával, és fordítottan arányos a távolságuk négyzetével:
\begin{equation}
    F = G \frac{M m}{r^2}
\end{equation}
Ahol $F$ a gravitációs erő, $G$ a gravitációs állandó, $M$ és $m$ a két test tömege, és $r$ a távolságuk. 

\subsubsection{Kepler-probléma megoldása}
A kepler-probléma lényegében a kéttest probléma, ahol két test egy olyan vonzó erő hatására mozog, amely a távolságuk négyzetével fordítottan arányos. A Kepler-probléma speciális esete, amikor az egyik test tömege sokkal nagyobb, mint a másiké (például egy bolygó és egy csillag esetében),
így a nehezebb test mozgása elhanyagolható. Ebben az esetben a könnyebb test mozgását a nehezebb test körül vizsgáljuk. A Kepler-probléma megoldását a Newtoni mechanika megalkotása tette lehetővé, amely a mozgás törvényein alapul.
A Newton gravitációs törvénye szerint a két test közötti gravitációs erő a következőképpen írható fel:
\begin{equation}
    \mathbf{F} = -G \frac{M m}{r^2} \hat{\mathbf{r}}
\end{equation}
Ahol $\mathbf{F}$ a gravitációs erő, $G$ a gravitációs állandó, $M$ és $m$ a két test tömege, és $\hat{\mathbf{r}}$ az egységvektor a két test közötti távolság irányában. 
A Newton második törvénye szerint a test mozgását a következőképpen írhatjuk fel:
\begin{equation}
    m \mathbf{a} = \mathbf{F}
\end{equation}
\begin{equation}
    \mathbf{a} = -G \frac{M}{r^2} \hat{\mathbf{r}}
\end{equation}
Mivel az erőtér centrális, ezért ezért tudjuk, hogy keringő mozgás csak akkor állhat fenn, ha a gyorsulás és a keringő tárgy sebessége minden pillanatban merőleges egymásra. Ekkor bevezethetünk egy vektort:
\begin{equation}
    \mathbf{K} = \mathbf{r} \times \mathbf{v}
\end{equation}
Ami az impulzusmomentummal arányos mennyiség. Erről a mennyiségről tudjuk, hogy időben állandó, hisz:

\begin{equation}
    \frac{d}{dt} \mathbf{K} = \frac{d}{dt} (\mathbf{r} \times \mathbf{v}) = \mathbf{v} \times \mathbf{v} + \mathbf{r} \times \mathbf{a} = \vec{0} \times \vec{0} = \vec{0}
\end{equation}
Illetve ha áttérünk a polár koordinátákra, akkor a helyvektor és a sebességvektor a következőképpen írható fel:
\begin{equation}
    \mathbf{r} = r \hat{\mathbf{e}}_r
\end{equation}
\begin{equation}
    \mathbf{v} = \dot{r} \hat{\mathbf{e}}_r + r \dot{\varphi} \hat{\mathbf{e}}_{\varphi}
\end{equation}
Ahol $\hat{\mathbf{e}}_r$ és $\hat{\mathbf{e}}_{\varphi}$ a polár koordináták egységvektorai. Ezt behelyettesítve a $\mathbf{K}$ vektorba:
\begin{equation}
    \mathbf{K} = r \hat{\mathbf{e}}_r \times (\dot{r} \hat{\mathbf{e}}_r + r \dot{\varphi} \hat{\mathbf{e}}_{\varphi}) = r^2 \dot{\varphi} (\hat{\mathbf{e}}_r \times \hat{\mathbf{e}}_{\varphi}) = r^2 \dot{\varphi} \hat{\mathbf{e}}_z
\end{equation}
Ahol $\hat{\mathbf{e}}_z$ a síkra merőleges egységvektor. Vagyis a $\mathbf{K}$ vektor nagysága:
\begin{equation}
    K = r^2 \dot{\varphi}
\end{equation}
Mivel a $\mathbf{K}$ vektor időben állandó, ezért a $K$ nagysága is állandó, így a következő összefüggést kapjuk:
\begin{equation}
    r^2 \dot{\varphi} = \text{állandó}
\end{equation}
Ez az összefüggés a területi sebesség állandóságát fejezi ki, ami a Kepler második törvényének matematikai megfogalmazása. Ez azt jelenti, hogy a bolygó és a Nap közötti képzeletbeli vonal egyenlő területeket súrol egyenlő idő alatt.

Ahhoz, hogy megkapjuk a pályaegyenletet használjuk ki a hármas skalárszorzat tulajdonságait:
\begin{equation}
    \mathbf{r} \times \mathbf{v} \cdot \mathbf{K} = \mathbf{v} \times \mathbf{K} \cdot \mathbf{r}
\end{equation}
Ekkor a baloldal:
\begin{equation}
    \mathbf{r} \times \mathbf{v} \cdot \mathbf{K} = \mathbf{K} \cdot \mathbf{K} = K^2
\end{equation}
A jobboldal kiértékeléséhez használjuk ki a $\mathbf{a} = -G \frac{M}{r^2} \hat{\mathbf{r}}$ gyorsulást, illetve azt hogy polárkoordinátákban $\frac{d}{dt}\hat{\mathbf{e}}_r = \frac{d\varphi}{dt}\hat{\mathbf{e}}_\varphi$ 
\begin{equation}
    \begin{aligned}
        \frac{d}{dt} (\mathbf{v} \times \mathbf{K}) &= \mathbf{a} \times \mathbf{K} = \left(- \frac{GM}{r^2} \hat{\mathbf{e}}_r\right) \times \left(r^2 \dot{\varphi} (\hat{\mathbf{e}}_r \times \hat{\mathbf{e}}_{\varphi})\right) = \\
        &= -GM \dot{\varphi} \left[\hat{\mathbf{e}}_r \times (\hat{\mathbf{e}}_r \times \hat{\mathbf{e}}_{\varphi})\right] = GM \dot{\varphi} \hat{\mathbf{e}}_{\varphi} = GM \frac{d}{dt} \hat{\mathbf{e}}_r
    \end{aligned}
\end{equation}
Ebből következik, hogy a $\frac{d}{dt} (\mathbf{v} \times \mathbf{K} - GM \hat{\mathbf{e}}_r) = \vec{0}$ vagyis ez a $\mathbf{v} \times \mathbf{K} - GM \hat{\mathbf{e}}_r$ egy konstans vektor,
amit nevezzünk el $GM\boldsymbol{\varepsilon}$:
\begin{equation}
    \mathbf{v} \times \mathbf{K} = GM (\boldsymbol{\varepsilon} + \hat{\mathbf{e}}_r)
\end{equation}
Tehát az eredeti hármas skalárszorzat egyenletünk jobboldala:
\begin{equation}
    \mathbf{v} \times \mathbf{K} \cdot \mathbf{r} = GM (\boldsymbol{\varepsilon} + \hat{\mathbf{e}}_r) \cdot \mathbf{r} = GM (|\mathbf{r}| + \boldsymbol{\varepsilon} \cdot \mathbf{r}))
\end{equation}
Egyenlővé téve a két oldalt:
\begin{equation}
    K^2 = GM (r + \boldsymbol{\varepsilon} \cdot \mathbf{r})
\end{equation}
Átrendezve:
\begin{equation}
    |\mathbf{r}| + \boldsymbol{\varepsilon} \cdot \mathbf{r} = \frac{K^2}{GM}
\end{equation}
Ahol $|\mathbf{r}| = r$ és $\boldsymbol{\varepsilon} \cdot \mathbf{r} = \varepsilon r cos\varphi$, így a pályaegyenlet:
\begin{equation}
    r (1 + \varepsilon cos\varphi) = \frac{K^2}{GM}
\end{equation}
Ez pedig egy kónikus metszet egyenlete, ahol $\varepsilon$ az excentricitás.

A kúpmetszeteknek viszont sok tipusa van, ezért, hogy pontosan meghatározzuk a pálya alakját, vizsgálódjunk egy kicsit tovább.
Ha az $x$ tengelyt úgy definiáljuk, hogy párhuzamos legyen a $\boldsymbol{\varepsilon} = \varepsilon \hat{\mathbf{e}}_x$ vektorral, illetve a konstansokat egy változóba foglaljuk:
\begin{equation}
    \gamma = \frac{K^2}{GM}
\end{equation}
Akkor a pályaegyenlet a következőképpen néz ki:
\begin{equation}
    r (1 + \varepsilon cos\varphi) = \gamma
\end{equation}
Ezt átírhatjuk derékszögű koordinátákba is:
\begin{equation}
    \sqrt{x^2 + y^2} (1 + \varepsilon \frac{x}{\sqrt{x^2 + y^2}}) = \gamma
\end{equation}
\begin{equation}
    \sqrt{x^2 + y^2} + \varepsilon x = \gamma
\end{equation}
\begin{equation}
    \sqrt{x^2 + y^2} = \gamma - \varepsilon x
\end{equation}
Négyzetre emelve:
\begin{equation}
    x^2 + y^2 = \gamma^2 - 2 \gamma \varepsilon x + \varepsilon^2 x^2
\end{equation}
\begin{equation}
    (1 - \varepsilon^2) x^2 + y^2 + 2 \gamma \varepsilon x - \gamma^2 = 0
\end{equation}
Ezt az egyenletet lehet tovább rendezni, hogy a kúpmetset típusát meghatározzuk:
\begin{equation}
    (1 - \varepsilon^2) (x^2 + \frac{2 \gamma \varepsilon}{1 - \varepsilon^2} x) + y^2 - \gamma^2 = 0
\end{equation}
\begin{equation}
    (1 - \varepsilon^2) \left(x^2 + \frac{2 \gamma \varepsilon}{1 - \varepsilon^2} x + \left(\frac{\gamma \varepsilon}{1 - \varepsilon^2}\right)^2\right) + y^2 - \gamma^2 - (1 - \varepsilon^2) \left(\frac{\gamma \varepsilon}{1 - \varepsilon^2}\right)^2 = 0
\end{equation}
\begin{equation}
    (1 - \varepsilon^2) \left(x + \frac{\gamma \varepsilon}{1 - \varepsilon^2}\right)^2 + y^2 = \gamma^2 + \frac{\gamma^2 \varepsilon^2}{1 - \varepsilon^2}
\end{equation}
Ahol bevezethezünk egy új változót az egyszerűsítés érdekében:
\begin{equation}
    \alpha = \frac{\gamma \varepsilon}{1 - \varepsilon^2}
\end{equation}
Így az egyenlet:
\begin{equation}
    (1 - \varepsilon^2) (x + \alpha)^2 + y^2 = \gamma^2 + \alpha \gamma \varepsilon
\end{equation}
A kúpmetszet egyenlete alapján a legnagyobb távolság a középponttól a következőképpen határozható meg:
\begin{equation}
    |\alpha| + \frac{\gamma}{1 - \varepsilon^2} = \frac{|\varepsilon| \gamma}{1 - \varepsilon^2} + \frac{\gamma}{1 - \varepsilon^2} = \frac{(1 + |\varepsilon|) \gamma}{(1 - |\varepsilon|)(1 + |\varepsilon|)} = \frac{\gamma}{1 - |\varepsilon|}
\end{equation}
A legkisebb távolság pedig:
\begin{equation}
    |\alpha| - \frac{\gamma}{1 - \varepsilon^2} = \frac{|\varepsilon| \gamma}{1 - \varepsilon^2} - \frac{\gamma}{1 - \varepsilon^2} = \frac{(|\varepsilon| - 1) \gamma}{(1 - |\varepsilon|)(1 + |\varepsilon|)} = -\frac{\gamma}{1 + |\varepsilon|}
\end{equation}
Az átlagos távolság lesz a féltengely:
\begin{equation}
    a = \frac{1}{2} \left(\frac{\gamma}{1 - |\varepsilon|} - \left(-\frac{\gamma}{1 + |\varepsilon|}\right)\right) = \frac{\gamma}{1 - \varepsilon^2}
\end{equation}
Most már meg tudjuk határozni a pálya típusát az excentricitás alapján:
\begin{itemize}
    \item Ha $0 \leq \varepsilon < 1$, akkor a pálya egy ellipszis, ami a bolygók és a Nap közötti pályára jellemző.
    \item Ha $\varepsilon = 1$, akkor a pálya egy parabola, ami olyan testekre jellemző, amelyek elég gyorsak ahhoz, hogy elhagyják a csillag gravitációs terét.
    \item Ha $\varepsilon > 1$, akkor a pálya egy hiperbola, ami olyan testekre jellemző, amelyek még gyorsabbak, és szintén elhagyják a csillag gravitációs terét.
\end{itemize}
\begin{figure}[H]
    \centering
    \includegraphics[width=0.8\textwidth]{imgs/3-tetel/kepler_palya.png}
    \caption{Különböző pályatípusok a Kepler-probléma esetében excentricitás függvényében}
    \label{fig:kepler_palya}
\end{figure}

\subsubsection*{Kepler Harmadik Törvénye}
Ebből a megoldásból könnyen megkapható Kepler harmadik törvénye is. Az ellipszis pálya területe:
\begin{equation}
    \pi \frac{\gamma}{1 - \varepsilon^2} \frac{\gamma}{\sqrt{1 - \varepsilon^2}} = \frac{\pi K^4}{G^2 M^2 (1 - \varepsilon^2)^{3/2}}
\end{equation}
Mivel a területi sebesség állandó és egyenlő $\frac{K}{2}$-val, ezért a keringési idő:
\begin{equation}
    T = \frac{2 \pi K^4}{G^2 M^2 (1 - \varepsilon^2)^{3/2}} \cdot \frac{2}{K} = \frac{4 \pi K^3}{G^2 M^2 (1 - \varepsilon^2)^{3/2}}
\end{equation}
Ebből következik, hogy:
\begin{equation}
    T^2 = \frac{16 \pi^2 K^6}{G^4 M^4 (1 - \varepsilon^2)^3}
\end{equation}
Mivel $a = \frac{K^2}{GM(1 - \varepsilon^2)}$, ezért $K^2 = a GM (1 - \varepsilon^2)$, így:
\begin{equation}
    T^2 = \frac{16 \pi^2 (a GM (1 - \varepsilon^2))^3}{G^4 M^4 (1 - \varepsilon^2)^3} = \frac{16 \pi^2 a^3 G^3 M^3 (1 - \varepsilon^2)^3}{G^4 M^4 (1 - \varepsilon^2)^3} = \frac{16 \pi^2 a^3}{G M}
\end{equation}
Ez pedig pontosan Kepler harmadik törvénye, amely kimondja, hogy a bolygók keringési idejének négyzete arányos a pálya fél nagytengelyének köbével.


\subsection{Kvantummechanikai problémák}
\subsubsection{Potenciálvölgy}
A potenciálvölgy általánosságban egy olyan fizikai rendszer, amelyben egy részecske egy adott potenciálfüggvény hatása alatt mozog, de a potenciálfüggvény egy adott tartományban a részecske energiája fölé emelkedik.
Ezt a régiót klasszikus fizikában potenciálgátként értelmezzük, amelyet a részecske nem képes átlépni, ha az energiája kisebb, mint a potenciálgát magassága. Azonban a kvantummechanikában a részecske hullámtermészete miatt létezik egy jelenség,
amelyet alagúthatásnak neveznek, amely lehetővé teszi a részecske számára, hogy áthaladjon a potenciálgáton még akkor is, ha az energiája kisebb, mint a gát magassága. Ez a jelenség a kvantummechanika egyik alapvető tulajdonsága, és számos fizikai jelenségben megfigyelhető,
például a radioaktív bomlásban, a szilárdtestfizikában és a kémiai reakciókban.
\begin{figure}[H]
    \centering
    \includegraphics[width=0.6\textwidth]{imgs/3-tetel/potencialvolgy.png}
    \caption{Potenciálvölgy Egy potenciálfüggvény esetében}
    \label{fig:potencialvolgy}
\end{figure}
\subsubsection*{Végtelen potenciálvölgy}
Nézzük meg először a végtelen potenciálvölgy esetét, ahol a potenciál-gát magasságát "effektíve" végtelennek tekintjük. Ez azt jelenti, hogy a részecske nem képes áthaladni a gátakon, és csak a potenciálvölgy belsejében mozoghat.
Bár ebben az esetben nincs alagúthatás, de a részecske hullámtermészete miatt így is váratlan eredményre fogunk jutni.
\begin{figure}[H]
    \centering
    \includegraphics[width=0.6\textwidth]{imgs/3-tetel/vegtelenpotencialvolgy.png}
    \caption{Végtelen potenciálvölgy}
    \label{fig:vegtelen_potencialvolgy}
\end{figure}
Ebben az esetben a potenciálfüggvény-t könnyen definiálhatjuk a következőképpen:
\begin{equation}
    V(x) = \begin{cases}
    0 & \text{ha } 0 < x < a \\
    \infty & \text{egyébként}
    \end{cases}
\end{equation}
Ahol $a$ a potenciálvölgy szélessége. A részecske mozgását az időfüggetlen Schrödinger-egyenlet írja le:
\begin{equation}
    -\frac{\hbar^2}{2m} \frac{d^2 \psi(x)}{dx^2} + V(x) \psi(x) = E \psi(x)
\end{equation}
Ahol $\psi(x)$ a részecske hullámfüggvénye, $E$ az energia, $m$ a részecske tömege, és $\hbar$ a redukált Planck-állandó. A falakon kívül $\psi (x) = 0$-nek kell lennie, tehát a részecskét nulla valószínűséggel találjuk meg ott.
A potenciálvölgy belsejében ($0 < x < a$) a Schrödinger-egyenlet a következőképpen egyszerűsödik:
\begin{equation}
    -\frac{\hbar^2}{2m} \frac{d^2 \psi(x)}{dx^2} = E \psi(x)
\end{equation}
\begin{equation}
    \frac{d^2 \psi(x)}{dx^2} + \frac{2mE}{\hbar^2} \psi(x) = 0
\end{equation}
Nevezzük át a paramétert:
\begin{equation}
    k \equiv \sqrt{\frac{2mE}{\hbar^2}}
\end{equation}
Itt feltételezzük, hogy az energia pozitív ($E > 0$), így $k$ valós szám lesz. Így az egyenlet a következőképpen néz ki:
\begin{equation}
    \frac{d^2 \psi(x)}{dx^2} + k^2 \psi(x) = 0
\end{equation}
A megoldás általános alakja:
\begin{equation}
    \psi(x) = A sin(kx) + B cos(kx)
\end{equation}
Ahol $A$ és $B$ integrálási állandók. Az állandókat a határfeltételek alapján tudjuk meghatározni. Általában azt feltételezzük, hogy a hullámfüggvénynek és a deriváltjának is folytonosnak kell lennie a határokon,
de a gát végtelensége miatt, most csak a hullámfüggvénytől követeljük meg a folytonosságot. Tehát a következő határfeltételeket alkalmazzuk:
\begin{equation}
    \psi(0) = \psi(a) = 0
\end{equation}
Az első határfeltétel alkalmazásával:
\begin{equation}
    \psi(0) = A sin(0) + B cos(0) = B = 0
\end{equation}
Tehát a hullámfüggvény egyszerűsödik:
\begin{equation}
    \psi(x) = A sin(kx)
\end{equation}
A második határfeltétel alkalmazásával:
\begin{equation}
    \psi(x) = A sin(kx) = 0
\end{equation}
Ezt teljesíthetné a $A = 0$, de ekkor csak annyit kaptunk, hogy $\psi (x) = 0$, ami triviális megoldás. teljesíthetné a $k = 0$ megoldás is, de ez feltételezné a $E = 0$-t, ami újból egy nem túl érdekes megoldás.
 A másik opció, hogy az x diszkrét értékeket vesz fel úgy, hogy a szinusz nulla legyen:
\begin{equation}
    k a = 0, \pm \pi, \pm 2\pi, \pm 3\pi, \ldots
\end{equation}
Itt az előjel nem számít, mivel a szinusz függvény páros. Tehát a diszkrét értékek:
\begin{equation}
    k_n = \frac{n \pi}{a} \quad \quad n = 1, 2, 3, \ldots
\end{equation}
Érdekesség, hogy a határfeltétel $x = a$ nem határozza meg az $A$ állandót, hanem a $k$-t adja meg. Ebből az eredményből pedig következik, hogy az energia csak bizonyos diszkrét értékeket vehet fel:
\begin{equation}
    E_n = \frac{\hbar^2 k_n^2}{2m} = \frac{\hbar^2 \pi^2 n^2}{2m a^2} \quad \quad n = 1, 2, 3, \ldots
\end{equation}
Ahhoz, hogy megkapjuk az $A$ állandót, normalizálnunk kell a hullámfüggvényt:
\begin{equation}
    \int_0^a |\psi_n(x)|^2 dx = 1
\end{equation}
\begin{equation}
    |A|^2 \int_0^a sin^2\left(\frac{n \pi x}{a}\right) dx = 1
\end{equation}
\begin{equation}
    |A|^2 \frac{a}{2} = 1
\end{equation}
\begin{equation}
    |A| = \sqrt{\frac{2}{a}}
\end{equation}
Így a normalizált hullámfüggvények:
\begin{equation}
    \psi_n(x) = \sqrt{\frac{2}{a}} sin\left(\frac{n \pi x}{a}\right) \quad \quad n = 1, 2, 3, \ldots
\end{equation}
Ezzel technikailag csak az $A$ nagyságát határoztuk meg, de mivel a hullámfüggvény csak a valószínűség sűrűséget határozza meg, azaz $|\psi (x)|^2$, ezért az $A$ előjelének nincs fizikai jelentősége.
Összefoglalva, a végtelen potenciálvölgy esetében a részecske energiája diszkrét értékeket vehet fel, amelyeket kvantált energiaszinteknek nevezünk. A hullámfüggvények pedig a potenciálvölgy belsejében oszcilláló
függvények, amelyek a határfeltételek miatt nullára mennek a gátakon kívül. Az $n = 1$ eset a legkisebb energiájú állapot, amelyet alapállapotnak nevezünk, míg az $n > 1$ esetek gerjesztett állapotoknak felelnek meg.
Ezeknek a $\psi_n (x)$ függvényeknek van egy pár érdekes tulajdonsága is:
\begin{itemize}
    \item Az állapotok ortogonálisak egymásra, azaz ha $n \neq m$, akkor:
    \begin{equation}
        \int_0^a \psi_n(x)^* \psi_m(x) dx = 0
    \end{equation}
    \item Ahogy megyünk feljebb az energiaszinteken (n növelése), a hullámfüggvények egyre több csomópontot tartalmaznak a potenciálvölgy belsejében. Az $n$-edik állapotban pontosan $(n-1)$ csomópont van. $n = 1$ esetében
    pedig nulla, mert a végpontok nem számítanak csomópontnak.
    \item a hullámfüggvények váltogatnak páros és páratlan függvények között, ahogy növeljük $n$ értékét. Az $n$ páratlan esetekben a hullámfüggvény páros, míg az $n$ páros esetekben páratlan függvény.
\end{itemize}
\begin{figure}[H]
    \centering
    \includegraphics[width=0.6\textwidth]{imgs/3-tetel/vegtelen_potencialvolgy_hullamfuggvenyek.png}
    \caption{Végtelen potenciálvölgy hullámfüggvényei az első néhány energiaszinten}
    \label{fig:vegtelen_potencialvolgy_hullamfuggvenyek}
\end{figure}

\subsubsection*{Véges potenciálvölgy}
A véges potenciálvölgy esetében a potenciál-gát magassága véges értékű, ami klasszikus fizikában nem jelentene nagy különbséget a végtelen potenciálvölgyhez képest, de a részecskék hullámtermészete miatt
érdekes jelenségeket tud produkálni a rendszer.

Vegyünk egy potenciálvölgyet amit a következőképpen definiálunk:
\begin{equation}
    V(x) = \begin{cases}
    -V_0 & - a < x < a \\
    0 & |x| \geq a
    \end{cases}
\end{equation}
\begin{figure}[H]
    \centering
    \includegraphics[width=0.4\textwidth]{imgs/3-tetel/vegespotencialvolgy.png}
    \caption{Véges potenciálvölgy}
    \label{fig:veges_potencialvolgy}
\end{figure}

Itt a potenciálfüggvény gátjának magasságát választjuk meg a nulla szintnek, és alatta $-V_0$-val bukik be a függvény $a$ és $-a$ között. Itt két lehetőséget határozhatunk meg:
\begin{itemize}
    \item Az energia $E$ kisebb, mint nulla ($E < 0$), ami azt jelenti, hogy a részecske kötött állapotban van a potenciálvölgyben.
    \item Az energia $E$ nagyobb, mint nulla ($E > 0$), ezt nevezik "szóródó állapotnak", ahol a részecske képes kilépni a potenciálvölgyből.
\end{itemize}
A részecske mozgását az időfüggetlen Schrödinger-egyenlet írja le:
\begin{equation}
    -\frac{\hbar^2}{2m} \frac{d^2 \psi(x)}{dx^2} + V(x) \psi(x) = E \psi(x)
\end{equation}
Ahol $\psi(x)$ a részecske hullámfüggvénye, $E$ az energia, $m$ a részecske tömege, és $\hbar$ a redukált Planck-állandó. A potenciálvölgy három régióra osztható:
\begin{itemize}
    \item Régió I: $x < -a$, ahol $V(x) = 0$
    \item Régió II: $-a < x < a$, ahol $V(x) = -V_0$
    \item Régió III: $x > a$, ahol $V(x) = 0$
\end{itemize}
Nézzük meg először a \textbf{kötött állapotot} ($E < 0$). A függvény az első régióban a következő egyenletet elégíti ki:
\begin{equation}
    \frac{d^2 \psi(x)}{dx^2} + \frac{2mE}{\hbar^2} \psi(x) = 0
\end{equation}
A megoldás a következő formát ölti:
\begin{equation}
    \psi(x) = A e^{- kx} + B e^{kx}
\end{equation}
Ahol $k = \sqrt{\frac{2mE}{\hbar^2}}$. Itt is beláthatjuk, hogy az $A$ együtthatós tag elszáll a végtelenben, így $A = 0$. Tehát az első régióban a hullámfüggvény:
\begin{equation}
    \psi_I(x) = B e^{kx}, \quad x < -a
\end{equation}
A harmadik régióban hasonlóan:
\begin{equation}
    \psi_{III}(x) = F e^{kx}, \quad x > a
\end{equation}
A második régióban a Schrödinger-egyenlet a következőképpen néz ki:
\begin{equation}
   -\frac{\hbar^2}{2m} \frac{d^2 \psi(x)}{dx^2} - V_0 \psi(x) = E \psi(x)
\end{equation}
\begin{equation}
    \frac{d^2 \psi(x)}{dx^2} + \frac{2m(E + V_0)}{\hbar^2} \psi(x) = 0
\end{equation}
Nevezzük át a paramétert:
\begin{equation}
    l \equiv \sqrt{\frac{2m(E + V_0)}{\hbar^2}}
\end{equation}
Itt megállapítjuk, hogy $E + V_0 > 0$, és bár $E < 0$, de nagyobbnak kell lennie $- V_0$-nál mert ha $E < V_{min}$ akkor a hullámfügvény nem lenne normálható, ami megszegi a hullámfüggvény valószínűség értelmezését. Így az egyenlet a következőképpen néz ki:
\begin{equation}
    \frac{d^2 \psi(x)}{dx^2} = - l^2 \psi(x)
\end{equation}
A megoldás általános alakja:
\begin{equation}
    \psi_{II}(x) = C sin(l x) + D cos(l x), \quad -a < x < a
\end{equation}
Itt élhetünk egy egyszerűsítéssel, mégpedig, hogy feltételezhetjük, hogy a hullámfüggvény vagy páros vagy páratlan lesz. Ezért, ha a potenciálvölgy szimmetrikus (és jelenleg úgy határoztuk meg),
akkor elég az egyik oldalra (pl. $+ a$-nál) megoldani as határfeltételeket, és abból automatikusan következik, hogy a másik oldalon $\psi(-x) = \pm \psi(x)$ lesz a megoldás (az előjel a paritás függvénye).
Nézzük meg a páros esetet, tehát vehetjük, hogy $C = 0$, így a teljes hullámfüggvény:
\begin{equation}
\psi(x) = \begin{cases}
    F e^{kx} & x > a \\
    D cos(l x) & 0 < x < a \\
    \psi_I(-x) & x < 0 
    \end{cases}
\end{equation}
Ezeknek az egyenleteknek a következő folytonossági határfeltételeket kell teljesíteniük az $x = \pm a$ pontokban:
\begin{equation}
    \psi_{II}(a) = \psi_{III}(a)
\end{equation}
\begin{equation}
    \psi_{II}'(a) = \psi_{III}'(a)
\end{equation}
Itt az elöbb tárgyalt párosság miatt elég, ha csak az $x = a$ pontban nézzük meg a feltételeket:
\begin{equation}
    D cos(l a) = F e^{- k a}
\end{equation}
\begin{equation}
    - D l sin(l a) = - k F e^{- k a}
\end{equation}
Ebből az $F$-et kifejezve és behelyettesítve a második egyenletbe:
\begin{equation}
    - D l sin(l a) = - k D cos(l a)
\end{equation}
\begin{equation}
    l \cdot tan(l a) = k
\end{equation}
Ez a formula mutatja a megengedett energiákat, hiszen $l$ és $k$ is az energia függvénye. Vezessünk be pár új változót a könnyebb ábrázolás érdekében:
\begin{equation}
    z \equiv l a = a \sqrt{\frac{2m(E + V_0)}{\hbar^2}}
\end{equation}
\begin{equation}
    z_0 \equiv a \sqrt{\frac{2m V_0}{\hbar^2}}
\end{equation}
Emellett a $k$ és az $l$ definíciója miatt a következő összefüggés is fennáll:
\begin{equation}
    k^2 + l^2 = \frac{2m V_0}{\hbar^2}
\end{equation}
\begin{equation}
    k^2 + \frac{z^2}{a^2} = \frac{z_0^2}{a^2}
\end{equation}
\begin{equation}
    ka = \sqrt{z_0^2 - z^2}
\end{equation}
Tehát a megengedett energiákra vonatkozó egyenlet a következő alakot ölti:
\begin{equation}
    z \cdot tan(z) = \sqrt{z_0^2 - z^2}
\end{equation}
Ez egy transzendens egyenlet, ahol $z_0$ a potenciálvölgy mélységét és szélességét határozza meg. Az egyenlet megoldásai a megengedett energiák diszkrét értékeit adják meg a kötött állapotok számára. 
Ezt a problémát csak numerikusan vagy grafikusan lehet megoldani, úgyhogy ábrázoljuk a bal és jobb oldalt külön-külön, és keressük meg azokat a pontokat, ahol a két görbe metszi egymást.
\begin{figure}[H]
    \centering
    \includegraphics[width=0.6\textwidth]{imgs/3-tetel/veges_potencialvolgy_megengedett_energiak.png}
    \caption{Véges potenciálvölgy megengedett energiák grafikus megoldása $z_0 = 8$ esetén}
\end{figure}
A megoldások közül kettő van amelyik különösen érdekes:
\begin{itemize}
    \item A potenciálvölgy \textbf{széles és mély}: Ha a potenciálvölgy nagy, akkor a metszéspontok egy kicsivel $z_n = \frac{n \pi}{2}$ fölött lesznek, ahol $n = 1. 3. 5 \ldots$ (a megoldás "másik fele" a páratlan függvény megoldásában van jelen).
    Innen belátható, hogy:
    \begin{equation}
        E_n + V_0 \approx \frac{\hbar^2 \pi^2 n^2}{2 m (2a)^2} \quad \quad n = 1, 3, 5, \ldots
    \end{equation}
    A völgy alja fölött, ami pont az mint a végtelen potenciálvölgy esetében, csak $2a$-val, vagyis pont a fele annak ($n$ páratlansága miatt csak az egyik felét kapjuk meg). Tehát a véges potenciálvölgy megoldása közeledik
    közeledik a végtelen potenciálvölgy megoldásához, ha $V_0 \to \infty$, de minden véges $V_0$ esetén véges számú kötött állapot van.
    \item A potenciálvölgy \textbf{vékony és sekély}: Ha a potenciálvölgy kicsi, akkor kevesebb és kevesebb kötött állapot lesz, míg el nem érjük a $z_0 < \frac{\pi}{2}$ értéket ahol csak egy marad. Viszont akármennyire kicsi is a potenciálvölgy,
    akkor is lesz legalább egy kötött állapot, tehát a részecske mindig képes lesz "bennragadni" a potenciálvölgyben.
\end{itemize}
A páratlan megoldáshoz nagyon hasonló eredményt kapunk:
\begin{equation}
    z \cdot cot(z) = - \sqrt{z_0^2 - z^2}
\end{equation}
Ahol a megengedett energiákra vonatkozó egyenlet a következő alakot ölti:
\begin{equation}
    E_n + V_0 \approx \frac{\hbar^2 \pi^2 n^2}{2 m (2a)^2} \quad \quad n = 2, 4, 6, \ldots
\end{equation}
Már csak annyi van hátra, hogy normalizáljuk a hullámfüggvényt, és meghatározzuk az együtthatókat. Ezt a következő integrállal tehetjük meg:
\begin{equation}
    \int_{-\infty}^{\infty} |\psi(x)|^2 dx = 1
\end{equation}
\begin{equation}
    \int_{-\infty}^{-a} |\psi_I(x)|^2 dx + \int_{-a}^{a} |\psi_{II}(x)|^2 dx + \int_{a}^{\infty} |\psi_{III}(x)|^2 dx = 1
\end{equation}
\begin{equation}
    |B|^2 \int_{-\infty}^{-a} e^{2 k x} dx + |D|^2 \int_{-a}^{a} cos^2(l x) dx + |F|^2 \int_{a}^{\infty} e^{-2 k x} dx = 1
\end{equation}
Itt $B$-t nem kell külön meghatározni, hiszen a paritás miatt $B = \pm F$ lesz, tehát a probléma tovább egyszerűsödik:
\begin{equation}
    |F|^2 \left( \int_{-\infty}^{-a} e^{2 k x} dx + \int_{a}^{\infty} e^{-2 k x} dx \right) + |D|^2 \int_{-a}^{a} cos^2(l x) dx = 1
\end{equation}
Ezek az integrálok könnyen kiszámíthatók:
\begin{equation}
    |F|^2 \left( \frac{e^{-2 k a}}{2 k} + \frac{e^{-2 k a}}{2 k} \right) + |D|^2 \left( a + \frac{sin(2 l a)}{2 l} \right) = 1
\end{equation}
\begin{equation}
    |F|^2 \frac{e^{-2 k a}}{k} + |D|^2 \left( a + \frac{sin(2 l a)}{2 l} \right) = 1
\end{equation}
A határfeltételekből pedig kifejezhetjük $F$-et $D$-ben:
\begin{equation}
    D cos(l a) = F e^{- k a}
\end{equation}
\begin{equation}
    F = D cos(l a) e^{k a}
\end{equation}
Behelyettesítve a normalizációs feltételbe:
\begin{equation}
    |D|^2 cos^2(l a) e^{2 k a} \frac{e^{-2 k a}}{k} + |D|^2 \left( a + \frac{sin(2 l a)}{2 l} \right) = 1
\end{equation}
\begin{equation}
    |D|^2 \left( \frac{cos^2(l a)}{k} + a + \frac{sin(2 l a)}{2 l} \right) = 1
\end{equation}
\begin{equation}
    |D| = \sqrt{\frac{1}{\frac{cos^2(l a)}{k} + a + \frac{sin(2 l a)}{2 l}}}
\end{equation}
\begin{equation}
    |F| = \sqrt{\frac{cos^2(l a) e^{2 k a}}{k \left( \frac{cos^2(l a)}{k} + a + \frac{sin(2 l a)}{2 l} \right)}}
\end{equation}
Ezzel megkaptuk a kötött állapotok hullámfüggvényeit és megengedett energiáit a véges potenciálvölgy esetében.
\newline
\newline
Most nézzük meg a \textbf{szóródás állapotot} ($E > 0$). Ebben az esetben is hasonló megoldásokat kapunk a három régióra:
\begin{equation}
\psi(x) = \begin{cases}
    A e^{i k x} + B e^{- i k x} & x < -a \\
    C sin(l x) + D cos(l x) & -a < x < a \\
    F e^{i k x} + G e^{- i k x} & x > a
    \end{cases}
\end{equation}
Ahol most is $k = \sqrt{\frac{2mE}{\hbar^2}}$ és $l = \sqrt{\frac{2m(E + V_0)}{\hbar^2}}$. A különbség annyi, hogy most
a falon kívül is oszcilláló megoldásokat kapunk, hiszen a részecske képes kilépni a potenciálvölgyből, így a hullámfüggvény nem megy nullára a falakon kívül.
Emellett feltételezzük, hogy a részecskenyaláb csak balról jön be, így $G = 0$. Ekkor a $x < -a$ régióban kapott hullámfüggvény
a bejövő és visszaverődő hullámokat írja le, míg a $x > a$ régióban kapott hullámfüggvény a kilépő hullámot írja le.
Tehát a hullámfüggvény most a következő alakot ölti:
\begin{equation}
\psi(x) = \begin{cases}
    A e^{i k x} + B e^{- i k x} & x < -a \\
    C sin(l x) + D cos(l x) & -a < x < a \\
    F e^{i k x} & x > a
    \end{cases}
\end{equation}
A határfeltételek alkalmazásával az $x = \pm a$ pontokban:
\begin{equation}
    \psi_{II}(-a) = \psi_{I}(-a)
\end{equation}
\begin{equation}
    \psi_{II}'(-a) = \psi_{I}'(-a)
\end{equation}
\begin{equation}
    \psi_{II}(a) = \psi_{III}(a)
\end{equation}
\begin{equation}
    \psi_{II}'(a) = \psi_{III}'(a)
\end{equation}
Tehát a folytonosság a $-a$ régióban:
\begin{equation}
    - C sin(l a) + D cos(l a) = A e^{- i k a} + B e^{i k a}
\end{equation}
A trigonometrikus függvények szimmetriája miatt lehet a negatív előjelet kivinni a szinusz elé, és eltüntetni a coszinuszban. A derivált folytonossága:
\begin{equation}
    ik \left[Ae^{- i k x} + B e^{i k x}\right] = l \left[C cos(l a) + D sin(l a)\right]
\end{equation}
A folytonosság az $+a$ régióban:
\begin{equation}
    F e^{i k a} = C sin(l a) + D cos(l a)
\end{equation}
\begin{equation}
    ik F e^{i k a} = l \left[C cos(l a) + D sin(l a)\right]
\end{equation}
Ezek segítségével eltüntethetjük $C$ és $D$-t, és kifejezhetjük $B$-t és $F$-et $A$-ban. Ebből pedig meghatározhatjuk a visszaverődési és átvitelési együtthatókat:
\begin{equation}
    B = i \frac{sin(2 l a)}{2 k l} \left(l^2 - k^2\right) F
\end{equation}
\begin{equation}
    F = \frac{e^{- i k a} A}{cos(2 l a) - i \frac{k^2 + l^2}{2 k l} sin(2 l a)}
\end{equation}
Ebből meghatározhatjuk a $T = \frac{|F|^2}{|A|^2}$ átviteli együtthatót és a $R = \frac{|B|^2}{|A|^2}$ visszaverődési együtthatót:
\begin{equation}
    T^{-1} = 1 + \frac{V_0^2}{4 E (E + V_0)} sin^2\left( \frac{2 a}{\hbar} \sqrt{2 m (E + V_0)} \right)
\end{equation}
Vegyük észre, hogy az átviteli együttható $T$ értéke $0$ és $1$ között van, hiszen a szinusz függvény négyzete sosem negatív.
és ha $T = 1$, vagyis a részecske teljesen áthalad a potenciálvölgyön anélkül, hogy visszaverődne, akkor a következő feltételnek kell teljesülnie:
\begin{equation}
    sin\left( \frac{2 a}{\hbar} \sqrt{2 m (E + V_0)} \right) = 0
\end{equation}
\begin{equation}
    \frac{2 a}{\hbar} \sqrt{2 m (E + V_0)} = n \pi \quad \quad n = 0, 1, 2, \ldots
\end{equation}
Ebből az energia értékek:
\begin{equation}
    E_n + V_0 = \frac{n^2 \pi^2 \hbar^2}{8 m a^2} \quad \quad n = 0, 1, 2, \ldots
\end{equation}
Ezek az energiák pontosan megegyeznek a végtelen potenciálvölgy esetében kapott energiákkal.
\begin{figure}[H]
    \centering
    \includegraphics[width=0.6\textwidth]{imgs/3-tetel/veges_potencialvolgy_atvitel.png}
    \caption{Véges potenciálvölgy átviteli együtthatója az energia függvényében}
    \label{fig:veges_potencialvolgy_atvitel_visszaverodes}
\end{figure}

\subsubsection{Oszcillátor}
A klasszikus mechanikai oszcillátort, ahol egy  $m$ tömeg egy  $k$ rugóállandóval rendelkező rugóhoz van kötve,
a Hooke-törvény alapján lehet felírni:
\begin{equation}
    F = - k x = m \frac{d^2 x}{dt^2}
\end{equation}
Ahol $x$ a tömeg kitérése az egyensúlyi helyzetből. A megoldás egy harmonikus mozgás lesz, ahol a tömeg $x(t) = A cos(\omega t + \phi)$ alakban mozog,
ahol $\omega = \sqrt{\frac{k}{m}}$ az oszcilláció körfrekvenciája, $A$ a amplitúdó és $\phi$ a kezdeti fázis. Az oszcillátor potenciálja a következőképpen írható fel:
\begin{equation}
    V(x) = \frac{1}{2} k x^2
\end{equation}
Ez természetesen csak egy "tökéletes" oszcillátorra igaz, ahol a rugó végtelenül rugalmas, és nincs súrlódás vagy egyéb veszteség a rendszerben.
A valóságban a Hooke-törvény gyakran már a rugó elszakadása előtt sem érvényes, ezért csak kis kitérésre lehet vele közelíteni. 
Ehhez írjuk fel a potenciál Taylor-sorát az egyensúlyi helyzet körül (V($x_0$) = 0):
\begin{equation}
    V(x) = V(x_0) + V'(x_0) (x - x_0) + \frac{1}{2} V''(x_0) (x - x_0)^2 + \ldots
\end{equation}
Itt az első tag az $V(x_0)$ defínicója miatt nulla, a magasabb rendű tagok pedig elhanyagolhatók kis kitérések esetén. Ezért a potenciál közelítőleg:
\begin{equation}
    V(x) \approx \frac{1}{2} V''(x_0) (x - x_0)^2
\end{equation}
Ami leírja a Harmonikus oszcillációt kis kitérések esetén, egy effektív rugóállandóval $k_{eff} = V''(x_0)$. Enek segítségével
virtuálisan bármilyen oszcilláló mozgást lehet harmonikus oszcillációval közelíteni, ha kicsi az amplitúdó. 
\newline
\newline
Ezek alapján egy kvantummechanikai harmonikus oszcillátort is közelítőleg felírhatunk a következő potenciállal:
\begin{equation}
    V(x) = \frac{1}{2} m \omega^2 x^2
\end{equation}
Ahol $\omega$ az oszcillátor körfrekvenciája és ezzel a nehezen definiálható kvantummechanikai rugóállandót "kivettük" a problémából.
A kvantummechanikai rendszerben nem a Newton-féle mozgásegyenletet használjuk, hanem az időfüggetlen Schrödinger-egyenletet:
\begin{equation}
    -\frac{\hbar^2}{2m} \frac{d^2 \psi(x)}{dx^2} + \frac{1}{2} m \omega^2 x^2 \psi(x) = E \psi(x)
\end{equation}
Ez egy másodrendű differenciálegyenlet, amelyet általában két módszerrrel szokás megoldani, a "brute force" hatványsor módszerrel,
ahol a hullámfüggvényt hatványsor formában írjuk fel, és a rekurzív együtthatókat határozzuk meg, vagy az elegánsabb algebrai módszerrel.
Jelen esetben nézzük meg az algebrai módszert, amely a létrehozó és megsemmisítő operátorok bevezetésén alapul.

\subsubsection*{Algebrai módszer}
Kezdjük azzal, hogy kicsit átrendezzük a Schrödinger-egyenletet:
\begin{equation}
    \frac{1}{2m} \left[\hat{p}^2 + (m \omega \hat{x})^2\right] \psi(x) = E \psi(x)
\end{equation}
Itt bevezettük az impulzus és hely operátort:
\begin{equation}
    \hat{x} = x
\end{equation}
\begin{equation}
    \hat{p} = - i \hbar \frac{d}{dx}
\end{equation}
Ez a Schrödinger-egyenlet alapján megfelel a Hamilton-operátornak:
\begin{equation}
    \hat{H}\psi = E\psi
\end{equation}
\begin{equation}
    \hat{H} = \frac{1}{2m} \left[\hat{p}^2 + (m \omega x)^2\right]
\end{equation}
Ahhoz, hogy meg tudjuk oldani az egyenletet normál esetben faktorizálnánk a Hamilton-operátort valahogy így:
\begin{equation}
    u^2 + v^2 = (iu + v)(- iu + v)
\end{equation}
Viszont itt nem egyszerű számokkal van dolgunk, mert $\hat{p}$ és $\hat{x}$ operátorok, és az operátoroknál nem feltétlenül igaz, hogy
kommutálnak egymással:
\begin{equation}
    \hat{p} \hat{x} \neq \hat{x} \hat{p} \qquad \text{általában}
\end{equation}
Ennek ellenére vezessünk be egy új operátor párt:
\begin{equation}
    \hat{a}_\pm = \frac{1}{\sqrt{2 \hbar m \omega}} \left( \mp i \hat{p} + m \omega \hat{x} \right)
\end{equation}
Nézzük meg, hogy mi a skaláris szorzatuk ezeknek az új operátoroknak:
\begin{equation}
    \begin{aligned}
    & \hat{a}_- \hat{a}_+ = \frac{1}{2 \hbar m \omega} \left( i \hat{p} + m \omega \hat{x} \right) \left( - i \hat{p} + m \omega \hat{x} \right) = \\
    & = \frac{1}{2 \hbar m \omega} \left( \hat{p}^2 + m^2 \omega^2 \hat{x}^2 + i m \omega (\hat{x} \hat{p} - \hat{p} \hat{x}) \right)
    \end{aligned}
\end{equation}
Itt látható, hogy megjelent egy extra $(\hat{x} \hat{p} - \hat{p} \hat{x})$ tag, ami kommutálás esetén eltűnne.
Ezt a tagot formálisan következőképpen írhatjuk fel:
\begin{equation}
    [\hat{x}, \hat{p}] = \hat{x} \hat{p} - \hat{p} \hat{x}
\end{equation}
És "kommutátornak" nevezzük. Ez a művelet azt méri, hogy két operátor mennyire "rosszul" kommutál egymással. Ha a kommutátor nulla, akkor a két operátor kommutál.
Most már felírhatjuk a bevezetett operátorok szorzatát a kommutátorral:
\begin{equation}
    \hat{a}_- \hat{a}_+ = \frac{1}{2 \hbar m \omega} \left( \hat{p}^2 + m^2 \omega^2 \hat{x}^2 + i m \omega [\hat{x}, \hat{p}] \right)
\end{equation}
\begin{equation}
    \hat{a}_- \hat{a}_+ = \frac{1}{\hbar \omega} \left[\hat{p}^2 + (m \omega \hat{x})^2\right] + \frac{i}{2 \hbar} [\hat{x}, \hat{p}]
\end{equation}
Itt pedig észrevehetjük a Hamilton-operátort:
\begin{equation}
    \hat{a}_- \hat{a}_+ = \frac{2}{\hbar \omega} \hat{H} + \frac{i}{2 \hbar} [\hat{x}, \hat{p}]
\end{equation}
Most már csak ki kell számolnunk a kommutátort. Ezt úgy tehetjük meg, hogy alkalmazzuk a két operátort egy tetszőleges teszt függvényre $f(x)$:
\begin{equation}
    [\hat{x}, \hat{p}] f(x) = \hat{x} \hat{p} f(x) - \hat{p} \hat{x} f(x)
\end{equation}
\begin{equation}
    = \hat{x} \left( - i \hbar \frac{d f(x)}{dx} \right) - \hat{p} \left( x f(x) \right)
\end{equation}
\begin{equation}
    = - i \hbar x \frac{d f(x)}{dx} + i \hbar \frac{d}{dx} \left( x f(x) \right)
\end{equation}
\begin{equation}
    = - i \hbar x \frac{d f(x)}{dx} + i \hbar \left( f(x) + x \frac{d f(x)}{dx} \right)
\end{equation}
\begin{equation}
    = i \hbar f(x)
\end{equation}
Mivel ez igaz minden $f(x)$ tesztfüggvényre, ezért a kommutátor maga:
\begin{equation}
    [\hat{x}, \hat{p}] = i \hbar
\end{equation}
Ezt visszahelyettesítve az előző egyenletbe:
\begin{equation}
    \hat{a}_- \hat{a}_+ = \frac{2}{\hbar \omega} \hat{H} + \frac{1}{2} [\hat{x}, \hat{p}] = \frac{2}{\hbar \omega} \hat{H} + \frac{1}{2} i \hbar
\end{equation}
vagy a Hamilton-operátorra rendezve:
\begin{equation}
    \hat{H} = \hbar \omega \left( \hat{a}_- \hat{a}_+ - \frac{1}{2} \right)
\end{equation}
Hasonlóan kiszámolhatjuk a másik sorrendű szorzatot is:
\begin{equation}
    \hat{a}_+ \hat{a}_- = \frac{1}{2 \hbar m \omega} \left( - i \hat{p} + m \omega \hat{x} \right) \left( i \hat{p} + m \omega \hat{x} \right) = 
\end{equation}
\begin{equation}
    = \frac{1}{2 \hbar m \omega} \left( \hat{p}^2 + m^2 \omega^2 \hat{x}^2 - i m \omega (\hat{x} \hat{p} - \hat{p} \hat{x}) \right)
\end{equation}
\begin{equation}
    \hat{a}_+ \hat{a}_- = \frac{1}{\hbar \omega} \left[\hat{p}^2 + (m \omega \hat{x})^2\right] - \frac{i}{2 \hbar} [\hat{x}, \hat{p}]
\end{equation}
\begin{equation}
    \hat{a}_+ \hat{a}_- = \frac{2}{\hbar \omega} \hat{H} - \frac{i}{2 \hbar} [\hat{x}, \hat{p}]
\end{equation}
\begin{equation}
    \hat{a}_+ \hat{a}_- = \frac{2}{\hbar \omega} \hat{H} - \frac{1}{2} i \hbar
\end{equation}
\begin{equation}
    \hat{H} = \hbar \omega \left( \hat{a}_+ \hat{a}_- + \frac{1}{2} \right)
\end{equation}

Látható, hogy a két eredmény pont 1-el tér el egymástól, ezért ezek az operátorok sem kommutálnak:
\begin{equation}
    [\hat{a}_-, \hat{a}_+] = 1
\end{equation}
Tehát összefoglalva a Hamilton-operátort a következőképpen írhatjuk fel:
\begin{equation}
    \hat{H} = \hbar \omega \left( \hat{a}_\pm \hat{a}_\mp \pm \frac{1}{2} \right)
\end{equation}
Ezzel a Schrödinger-egyenlet a következő alakot ölti:
\begin{equation}
    \hbar \omega \left( \hat{a}_\pm \hat{a}_\mp \pm \frac{1}{2} \right) \psi = E \psi
\end{equation}
Ekkor azt állítom, hogy ha $\psi$ kielégíti a Schrödinger-egyenletet egy adott $E$ energiára (tehát $\hat{H} \psi = E \psi$), akkor
$\hat{a}_+\psi$ is kielégíti a Schrödinger-egyenletet egy $E + \hbar \omega$ energiára: $(E + \hbar \omega)(\hat{a}_+\psi)$
\newline
\newline
Bizonyítás:
\begin{equation}
    \begin{aligned}
    & \hat{H} (\hat{a}_+ \psi) = \hbar \omega \left( \hat{a}_+ \hat{a}_- + \frac{1}{2} \right) (\hat{a}_+ \psi) = \hbar \omega \left( \hat{a}_+ \hat{a}_- \hat{a}_+ + \frac{1}{2} \hat{a}_+  \right)\psi \\
    & = \hbar \omega \hat{a}_+ \left(\hat{a}_- \hat{a}_+ + \frac{1}{2}\right)\psi = \hat{a}_+ \left[\hbar \omega \left(\hat{a}_- \hat{a}_+ + 1 + \frac{1}{2}\right)\psi\right] = \quad \quad \text{mert } \hat{a}_+ \hat{a}_- = \hat{a}_- \hat{a}_+ + 1 \\
    & = \hat{a}_+ \left( \hat{H} + \hbar \omega \right) \psi = \hat{a}_+ (E + \hbar \omega) \psi = \\
    & = (E + \hbar \omega)(\hat{a}_+ \psi)
    \end{aligned}
\end{equation}
Fontos megjegyezni, hogy a $\hat{a}_\pm$ operátorok szorzatának sorrendje számít, hiszen nem kommutálnak egymással, de az operátor és egy állandó 
szorzatánál nem, mert minden operátor kommutál minden állandóval.

Ugyanezt megtehetjük a $\hat{a}_-$ operátorral is:
\begin{equation}
    \begin{aligned}
    & \hat{H} (\hat{a}_- \psi) = \hbar \omega \left( \hat{a}_- \hat{a}_+ - \frac{1}{2} \right) (\hat{a}_- \psi) = \hbar \omega \left( \hat{a}_- \hat{a}_+ \hat{a}_- - \frac{1}{2} \hat{a}_-  \right)\psi \\
    & = \hbar \omega \hat{a}_- \left(\hat{a}_- \hat{a}_+ + \frac{1}{2}\right)\psi = \hat{a}_- \left[\hbar \omega \left(\hat{a}_- \hat{a}_+ - 1 - \frac{1}{2}\right)\psi\right] = \\
    & = \hat{a}_- \left( \hat{H} - \hbar \omega \right) \psi = \hat{a}_- (E - \hbar \omega) \psi = \\
    & = (E - \hbar \omega)(\hat{a}_- \psi)
    \end{aligned}
\end{equation}
Ezzel sikerült előállítani egy "gépezetet" amivel egy kezdő energiából tetszőleges magasabb vagy alacsonyabb energiájú állapotokat tudunk előállítani.
Ezért ezeket az operátorokat "létra operátoroknak" is szokták nevezni és a $\hat{a}_+$-at keltő és a $\hat{a}_-$-at eltüntető operátornak is hívják.
\begin{figure}[H]
    \centering
    \includegraphics[width=0.4\textwidth]{imgs/3-tetel/harmonikus_oszcillator_letra_operatorok.png}
    \caption{Harmonikus oszcillátor létra operátorok ábrázolása}
\end{figure}

De mi történik, ha ezeket az operátorokat végtelenszer alkalmazzuk? A keltő operáttorral nem lesz gond, hiszen azzal csak egyre nagyonbb és nagyobb
energiájú állapotokat állíthatunk elő. Viszont az eltüntető operátorral egy idő után el fogunk jutni egy olyan állapothoz, ahol az energia negatív lesz,
ami fizikailag értelmetlen. Ezért kell lennie egy olyan állapotnak amire ha alkalmazzuk az eltüntető operátort, akkor nullát kapunk vissza:
\begin{equation}
    \hat{a}_- \psi_0 = 0
\end{equation}
Ezt az állapotot nevezik alapállapotnak. Ahhoz, hogy meghatározzuk az alapállapot energiáját, alkalmazzuk a Hamilton-operátort az alapállapotra:
\begin{equation}
    \hat{H} \psi_0 = \hbar \omega \left( \hat{a}_+ \hat{a}_- + \frac{1}{2} \right) \psi_0 = \hbar \omega \left( 0 + \frac{1}{2} \right) \psi_0
\end{equation}
\begin{equation}
    \hat{H} \psi_0 = \frac{1}{2} \hbar \omega \psi_0
\end{equation}
Tehát az alapállapot energiája:
\begin{equation}
    E_0 = \frac{1}{2} \hbar \omega
\end{equation}
Most már csak az alapállapot hullámfüggvényét kell meghatároznunk. Ehhez használjuk fel az alapállapotra vonatkozó feltételt:
\begin{equation}
    \hat{a}_- \psi_0 = 0
\end{equation}
\begin{equation}
    \frac{1}{\sqrt{2 \hbar m \omega}} \left( i \hbar \frac{d}{dx} + m \omega x \right) \psi_0 = 0
\end{equation}
\begin{equation}
    \left( \frac{d}{dx} + \frac{m \omega}{i \hbar} x \right) \psi_0 = 0
\end{equation}
Ez egy elsőrendű differenciálegyenlet, amit könnyen meg tudunk oldani:
\begin{equation}
    \frac{d \psi_0}{\psi_0} = - \frac{m \omega}{i \hbar} x dx
\end{equation}
\begin{equation}
    \int \frac{d \psi_0}{\psi_0} = - \frac{m \omega}{i \hbar} \int x dx
\end{equation}
\begin{equation}
    ln(\psi_0) = - \frac{m \omega}{2 \hbar} x^2 + C
\end{equation}
\begin{equation}
    \psi_0(x) = A e^{- \frac{m \omega}{2 \hbar} x^2}
\end{equation}
És ha normáljuk a hullámfüggvényt:
\begin{equation}
    \int_{-\infty}^{\infty} |\psi_0(x)|^2 dx = 1
\end{equation}
\begin{equation}
    |A|^2 \int_{-\infty}^{\infty} e^{- \frac{m \omega}{\hbar} x^2} dx = 1
\end{equation}
\begin{equation}
    |A|^2 \sqrt{\frac{\pi \hbar}{m \omega}} = 1
\end{equation}
\begin{equation}
    |A| = \left( \frac{m \omega}{\pi \hbar} \right)^{1/4}
\end{equation}
Tehát az alapállapot hullámfüggvénye:
\begin{equation}
    \psi_0(x) = \left( \frac{m \omega}{\pi \hbar} \right)^{1/4} e^{- \frac{m \omega}{2 \hbar} x^2}
\end{equation}
Tehát az összes állapot előállítható az alapállapotból a keltő operátor ismételt alkalmazásával:
\begin{equation}
    \psi_n = A_n (\hat{a}_+)^n \psi_0 \quad \quad \text{ahol }  E_n = \hbar \omega \left( n + \frac{1}{2} \right) \quad \quad n = 0, 1, 2, \ldots
\end{equation}
Itt $A_n$ egy normalizáló faktor, amit a hullámfüggvény normálásával lehet meghatározni. Tehát például az első izgatott állapot hullámfüggvénye:
\begin{equation}
    \psi_1 = A_1 \hat{a}_+ \psi_0
\end{equation}
\begin{equation}
    = A_1 \frac{1}{\sqrt{2 \hbar m \omega}} \left( - i \hbar \frac{d}{dx} + m \omega x \right) \left( \frac{m \omega}{\pi \hbar} \right)^{1/4} e^{- \frac{m \omega}{2 \hbar} x^2}
\end{equation}
\begin{equation}
    = A_1 \frac{1}{\sqrt{2 \hbar m \omega}} \left( m \omega x + i \hbar \frac{m \omega}{\hbar} x \right) \left( \frac{m \omega}{\pi \hbar} \right)^{1/4} e^{- \frac{m \omega}{2 \hbar} x^2}
\end{equation}
\begin{equation}
    = A_1 \sqrt{\frac{2 m \omega}{\hbar}} x \left( \frac{m \omega}{\pi \hbar} \right)^{1/4} e^{- \frac{m \omega}{2 \hbar} x^2}
\end{equation}
Normálva ezt a hullámfüggvényt megkapjuk az első izgatott állapot a normalizáló faktort:
\begin{equation}
    \int |\psi_1(x)|^2 dx = |A_1|^2 \frac{2 m \omega}{\hbar} \left( \frac{m \omega}{\pi \hbar} \right)^{1/2} \int_{-\infty}^{\infty} x^2 e^{- \frac{m \omega}{\hbar} x^2} dx = 1
\end{equation}
\begin{equation}
    |A_1|^2 \frac{2 m \omega}{\hbar} \left( \frac{m \omega}{\pi \hbar} \right)^{1/2} \frac{\sqrt{\pi}}{2} \left( \frac{\hbar}{m \omega} \right)^{3/2} = 1
\end{equation}
\begin{equation}
    |A_1|^2 = 1
\end{equation}
Ahhoz, hogy általánosan normalizálni tudjuk a hullámfüggvényt, egy kicsit trükközni kell. Először is megállapíthatjuk, hogy
ha "feljebb" vagy "lejjebb" lépünk az  energia-létrán, akkor a hullámfüggvények arányosak a következő módon:
\begin{equation}
    \hat{a}_+\psi_{n} = c_n \psi_{n+1} \quad \quad \hat{a}_- \psi_{n} = d_n \psi_{n-1}
\end{equation}
Ahol $c_n$ és $d_n$ arányossági tényezők. Ezek meghatározásához először megjegyezzük, hogy tetszőleges függvények esetén igaz, hogy:
\begin{equation}
    \int_{-\infty}^{\infty} f^*(x) (\hat{a}_\pm g(x)) dx = \int_{-\infty}^{\infty} (\hat{a}_\mp f(x))^* g(x) dx
\end{equation}
Ahol $f^*(x)$ a komplex konjugáltja az $f(x)$ függvénynek. Ezt algebrában úgy nevezik, hogy $\hat{a}_\pm$ a hermitikus konjugáltja $\hat{a}_\mp$-nek.
Innen következik, hogy:
\begin{equation}
    \int_{-\infty}^{\infty} (\hat{a}_\pm \psi_n)^* (\hat{a}_\pm \psi_n) dx = \int_{-\infty}^{\infty} (\hat{a}_\mp \hat{a}_\pm \psi_n)^* \psi_n dx
\end{equation}
De mivel $\hat{a}_\mp \hat{a}_\pm = \frac{1}{\hbar \omega} \hat{H} \pm \frac{1}{2}$, ezért:
\begin{equation}
    \int_{-\infty}^{\infty} (\hat{a}_\pm \psi_n)^* (\hat{a}_\pm \psi_n) dx = \int_{-\infty}^{\infty} \left( \frac{E_n}{\hbar \omega} \pm \frac{1}{2} \right) \psi_n^* \psi_n dx
\end{equation}
\begin{equation}
    = \left( n + \frac{1}{2} \pm \frac{1}{2} \right) \int_{-\infty}^{\infty} \psi_n^* \psi_n dx
\end{equation}
\begin{equation}
    = n + \frac{1 \pm 1}{2}
\end{equation}
Vagyis:
\begin{equation}
    \hat{a}_+ \hat{a}_- \psi_n = n \psi_n \quad \quad \hat{a}_- \hat{a}_+ \psi_n = (n + 1) \psi_n
\end{equation}
Ezt behelyettesítve az előző egyenletbe:
\begin{equation}
    \int_{-\infty}^{\infty} (\hat{a}_+ \psi_n)^* (\hat{a}_+ \psi_n) dx = |c_n|^2 \int_{-\infty}^{\infty} |\psi_{n+1}|^2 dx = (n + 1) \int_{-\infty}^{\infty} |\psi_n|^2 dx
\end{equation}
\begin{equation}
    \int_{-\infty}^{\infty} (\hat{a}_- \psi_n)^* (\hat{a}_- \psi_n) dx = |d_n|^2 \int_{-\infty}^{\infty} |\psi_{n-1}|^2 dx = n \int_{-\infty}^{\infty} |\psi_n|^2 dx
\end{equation}
De mivel a hullámfüggvények normálva vannak, ezért az integrálok értéke 1, így:
\begin{equation}
    |c_n|^2 = n + 1 \quad \quad |d_n|^2 = n
\end{equation}
\begin{equation}
    \hat{a}_+ \psi_n = \sqrt{n + 1} \psi_{n+1} \quad \quad \hat{a}_- \psi_n = \sqrt{n} \psi_{n-1}
\end{equation}
Innen ha megnézzük az egyes állapotokra:
\begin{equation}
    \psi_1 = \hat{a}_+ \psi_0 = \sqrt{1} \psi_1 \implies \psi_1 = \hat{a}_+ \psi_0
\end{equation}
\begin{equation}
    \psi_2 = \frac{1}{\sqrt{2}} \hat{a}_+ \psi_1 = \frac{1}{\sqrt{2}} \hat{a}_+^2 \psi_0
\end{equation}
\begin{equation}
    \psi_3 = \frac{1}{\sqrt{3}} \hat{a}_+ \psi_2 = \frac{1}{\sqrt{3 \cdot 2}} \hat{a}_+^3 \psi_0
\end{equation}
\begin{equation}
    \psi_4 = \frac{1}{\sqrt{4}} \hat{a}_+ \psi_3 = \frac{1}{\sqrt{4 \cdot 3 \cdot 2}} \hat{a}_+^4 \psi_0
\end{equation}
És így tovább. Ebből következik az általános formula:
\begin{equation}
    \psi_n = \frac{1}{\sqrt{n!}} (\hat{a}_+)^n \psi_0 \quad \quad n = 0, 1, 2, \ldots
\end{equation}
Ezzel megkaptuk a harmonikus oszcillátor összes energiáját és hullámfüggvényét az algebrai módszer segítségével. A harmonikus oszcillátor
az egyik legfontosabb modell a kvantummechanikában, hiszen sok rendszer közelíthető vele, és a kvantumtérelmélet alapját is képezi. A segítségével
lehet közelíteni például a molekulákban elhelyezkedő atomok rezgéseit, vagy a kristályrácsok fononjait is.

\subsubsection*{Hatványsoros módszer}
Ezt a módszert csak nagyvonalakban tárgyaljuk mert ugyan arra az eredményre jutunk mint az előző módszerrel. Előnye,
hogy ez egy "általánosabb" módszer, amit más potenciálokra is lehet alkalmazni, míg az előző módszer csak a harmonikus oszcillátorra alkalmazható.
Kezdjük azzal, hogy írjuk át a Schrödinger-egyenletet egy kicsit:
\begin{equation}
    \frac{d^2 \psi(x)}{dx^2} = \frac{2m}{\hbar^2} \left( \frac{1}{2} m \omega^2 x^2 - E \right) \psi(x)
\end{equation}
\begin{equation}
    \frac{d^2 \psi(x)}{dx^2} = \left( \frac{m^2 \omega^2}{\hbar^2} x^2 - \frac{2mE}{\hbar^2} \right) \psi(x)
\end{equation}
Most vezessük be a következő változókat:
\begin{equation}
    \xi = \sqrt{\frac{m \omega}{\hbar}} x
\end{equation}
\begin{equation}
    \epsilon = \frac{2E}{\hbar \omega}
\end{equation}
Ezekkel a változókkal az egyenlet a következő alakot ölti:
\begin{equation}
    \frac{d^2 \psi(\xi)}{d \xi^2} = (\xi^2 - \epsilon) \psi(\xi)
\end{equation}
Most nézzük meg az egyenlet viselkedését nagy $\xi$ értékeknél. Ilyenkor a $\xi^2$ tag dominál, így az egyenlet közelítőleg:
\begin{equation}
    \frac{d^2 \psi(\xi)}{d \xi^2} \approx \xi^2 \psi(\xi)
\end{equation}
Ennek az egyenletnek a megoldása a következő alakú:
\begin{equation}
    \psi(\xi) \approx A e^{- \frac{\xi^2}{2}} + B e^{\frac{\xi^2}{2}}
\end{equation}
De mivel a hullámfüggvénynek végesnek kell lennie, ezért a $B$ együtthatót nullának kell vennünk, így:
\begin{equation}
    \psi(\xi) \approx A e^{- \frac{\xi^2}{2}}
\end{equation}
Ezért a teljes hullámfüggvényt felírhatjuk a következő alakban:
\begin{equation}
    \psi(\xi) = h(\xi) e^{- \frac{\xi^2}{2}}
\end{equation}
Ahol $h(\xi)$ egy ismeretlen függvény, amit meg kell határoznunk. Ezt visszahelyettesítve az eredeti egyenletbe:
\begin{equation}
    \frac{d^2 h(\xi)}{d \xi^2} - 2 \xi \frac{d h(\xi)}{d \xi} + (\epsilon - 1) h(\xi) = 0
\end{equation}
Most feltételezzük, hogy $h(\xi)$ egy hatványsor alakban írható fel:
\begin{equation}
    h(\xi) = \sum_{n=0}^{\infty} a_n \xi^n
\end{equation}
Ezt behelyettesítve az egyenletbe és rendezzük az egyenletet:
\begin{equation}
    \sum_{n=0}^{\infty} \left[ (n+2)(n+1) a_{n+2} - 2 n a_n + (\epsilon - 1) a_n \right] \xi^n = 0
\end{equation}
Ez akkor igaz, ha minden együttható nulla:
\begin{equation}
    (n+2)(n+1) a_{n+2} - 2 n a_n + (\epsilon - 1) a_n = 0
\end{equation}
Ebből kifejezve az $a_{n+2}$ együtthatót:
\begin{equation}
    a_{n+2} = \frac{2 n + 1 - \epsilon}{(n+2)(n+1)} a_n
\end{equation}
Most nézzük meg a sorozat viselkedését nagy $n$ értékeknél. Ilyenkor az együttható közelítőleg:
\begin{equation}
    a_{n+2} \approx \frac{2 n + 1 - \epsilon}{(n+2)(n+1)} a_n
\end{equation}
\begin{equation}
    \approx \frac{2}{n} a_n
\end{equation}
Ez azt jelenti, hogy az együtthatók növekedése olyan gyors, hogy a sorozat divergens lesz, és a hullámfüggvény nem lesz normálható. Ezért a sorozatot le kell vágni egy bizonyos $n$ értéknél, ami csak akkor lehetséges, ha a számláló nulla lesz egy adott $n$ értéknél:
\begin{equation}
    2 n + 1 - \epsilon = 0
\end{equation}
\begin{equation}
    \epsilon = 2 n + 1
\end{equation}
Ebből következik az energia kvantálódása:
\begin{equation}
    E_n = \hbar \omega \left( n + \frac{1}{2} \right) \quad \quad n = 0, 1, 2, \ldots
\end{equation}
Ez megegyezik az előző módszerrel kapott eredménnyel. A hullámfüggvényeket pedig a hatványsorozatból lehet meghatározni az együtthatók segítségével:
\begin{equation}
    \psi_n(\xi) = h_n(\xi) e^{- \frac{\xi^2}{2}} \quad \quad \text{ahol } h_n(\xi) \text{ egy n-ed fokú polinom}
\end{equation}
Ezek a polinomok a Hermite-polynomok, és a hullámfüggvények a következő alakot öltik:
\begin{equation}
    \psi_n(x) = \left( \frac{m \omega}{\pi \hbar} \right)^{1/4} \frac{1}{\sqrt{2^n n!}} H_n\left( \sqrt{\frac{m \omega}{\hbar}} x \right) e^{- \frac{m \omega}{2 \hbar} x^2}
\end{equation}
Ahol $H_n(\xi)$ az n-ed fokú Hermite-polinom. Így tehát a hatványsoros módszerrel is megkaptuk a harmonikus oszcillátor energiáit és hullámfüggvényeit.
Néhány első Hermite-polinom:
\begin{equation*}
    \begin{aligned}
    & H_0(\xi) = 1 \\
    & H_1(\xi) = 2 \xi \\
    & H_2(\xi) = 4 \xi^2 - 2 \\
    & H_3(\xi) = 8 \xi^3 - 12 \xi \\
    & H_4(\xi) = 16 \xi^4 - 48 \xi^2 + 12 \\
    & H_5(\xi) = 32 \xi^5 - 160 \xi^3 + 120 \xi \\
    \end{aligned}    
\end{equation*}

\subsubsection{Rotátor}
Képzeljünk el egy olyan kvantummechanikai rendszert, ahol két részecske ($m_1$ és $m_2$ tömeggel) egy tömeg nélküli, $a$ hosszúságú rúddal van összekötve,
és szabadon foroghat a tér három irányába a kötött súlypontja körül. Ezt a rendszert nevezik \textbf{rotátornak}, és egyszerű molekulál modellezésére gyakran használatos.
\begin{figure}[H]
    \centering
    \includegraphics[width=0.8\textwidth]{imgs/3-tetel/rotator.png}
    \caption{Rotátor modell}
\end{figure}

Először nézzük meg a megengedett energia szinteket. A klasszikus rotátor kinetikus energiája a következőképpen írható fel:
\begin{equation}
    E_K = \frac{L^2}{2I}
\end{equation}
Ahol $L$ a rotátor perdülete, és $I$ a tehetetlenségi nyomatéka, ami a következőképpen számítható ki:
\begin{equation}
    I = \mu a^2
\end{equation}
Ahol $\mu$ a két részecske redukált tömege:
\begin{equation}
    \mu = \frac{m_1 m_2}{m_1 + m_2}
\end{equation}
Az itt megjelenő $L^2$ operátor szerencsére eléggé ismerős a kvantummechanikában, hiszen ezt az operátort gyakran használjuk a szögmomentum
vizsgálatakor. Ennek az operátornak az sajátértéke:
\begin{equation}
    L^2 = \hbar^2 l (l + 1) \quad \quad l = 0, 1, 2, \ldots, n-1
\end{equation}
Ahol $l$ a kvantumszám, ami az impulzusmomentumhoz tartozik. Ezt behelyettesítve a kinetikus energiára:
\begin{equation}
    E_l = \frac{\hbar^2}{2I} l (l + 1)
\end{equation}
Mivel az energiaszinteket általában $n$-nel jelöljük, ezért átírhatjuk az egyenletet a következő alakra:
\begin{equation}
    E_n = \frac{\hbar^2}{2I} n (n + 1) \quad \quad n = 0, 1, 2, \ldots \quad \text{ahol} \quad I = \frac{m_1 + m_2}{m_1 m_2} a^2
\end{equation}

Most hogy a sajátértékeket ismerjük, nézzük meg a sajátfüggvényeket is. Karakterizálják az irányt a térben a $\theta$ és $\phi$ szögek.
Itt is kihasználhatjuk, hogy a rotátor energiája csak a $L^2$ operátortól függ, így a sajátfüggvény egyenlet:
\begin{equation}
    L^2 = - \hbar^2 \left[ \frac{1}{\sin \theta} \frac{\partial}{\partial \theta} \left( \sin \theta \frac{\partial}{\partial \theta} \right) + \frac{1}{\sin^2 \theta} \frac{\partial^2}{\partial \phi^2} \right]
\end{equation}
Tehát a sajátérték egyenlet a következő alakot ölti:
\begin{equation}
    - \frac{\hbar^2}{2I} \left[ \frac{1}{\sin \theta} \frac{\partial}{\partial \theta} \left( \sin \theta \frac{\partial Y(\theta, \phi)}{\partial \theta} \right) + \frac{1}{\sin^2 \theta} \frac{\partial^2 Y(\theta, \phi)}{\partial \phi^2} \right] = E Y(\theta, \phi)
\end{equation}
De az energiaszinteket már ismerjük, így behelyettesítve:
\begin{equation}
    \left[ \frac{1}{\sin \theta} \frac{\partial}{\partial \theta} \left( \sin \theta \frac{\partial Y(\theta, \phi)}{\partial \theta} \right) + \frac{1}{\sin^2 \theta} \frac{\partial^2 Y(\theta, \phi)}{\partial \phi^2} \right] + l (l + 1) Y(\theta, \phi) = 0
\end{equation}
Ez az egyenlet ismerős, hiszen ez a szögfüggvényekre vonatkozó sajátérték egyenlet, aminek a megoldása a szögfüggvények:
\begin{equation}
    Y_{l}^{m}(\theta, \phi) = \sqrt{\frac{(2l + 1)}{4 \pi} \frac{(l - m)!}{(l + m)!}} P_{l}^{m}(\cos \theta) e^{i m \phi}
\end{equation}
Ahol $P_{l}^{m}(\cos \theta)$ a kapcsolt Legendre-polinom, és $m$ az impulzusmomentum mágneses kvantumszáma, ami a következő értékeket veheti fel:
\begin{equation}
    m = -l, -l + 1, \ldots, 0, \ldots, l - 1, l
\end{equation}
Összefoglalva tehát a rotátor energiái és hullámfüggvényei a következőképpen alakulnak:
\begin{equation}
    E_l = \frac{\hbar^2}{2I} l (l + 1) \quad \quad l = 0, 1, 2, \ldots
\end{equation}
\begin{equation}
    Y_{l}^{m}(\theta, \phi) = \sqrt{\frac{(2l + 1)}{4 \pi} \frac{(l - m)!}{(l + m)!}} P_{l}^{m}(\cos \theta) e^{i m \phi} \quad \quad m = -l, -l + 1, \ldots, l
\end{equation}
A rendszer degenerációja pedig $2l + 1$, hiszen minden $l$ értékhez $2l + 1$ különböző $m$ érték tartozik.

Mivel ismerjük az egyes energiaszinteket, ezért a rendszer spektrumát is meg tudjuk határozni. A spektrumot az adja, hogy a rendszer elnyel vagy kibocsájt egy fotont, aminek lesz egy meghatározott energiája és ezáltal hullámhossza.
Egy foton energiája a Planck-féle reláció alapján:
\begin{equation}
    E = h \nu = \frac{hc}{\lambda}
\end{equation}
Ahol $\nu$ a foton frekvenciája, $c$ a fény sebessége, és $\lambda$ a foton hullámhossza. Ha a rotátor egyik energiaszintjéről egy másikra lép, akkor a foton energiája a két energiaszint különbsége lesz:
\begin{equation}
    \Delta E = E_{l + 1} - E_{l} = \frac{\hbar^2}{2I} \left[ (l+1) (l + 2) - l (l + 1) \right] = \frac{\hbar^2}{2I} (2l + 2) = \frac{\hbar^2}{I} (l + 1)
\end{equation}
Ebből következik a foton frekvenciája:
\begin{equation}
    \nu_l = \frac{\Delta E}{h} = \frac{\hbar^2}{h I} (l + 1) = \frac{\hbar}{2 \pi I} (l + 1)
\end{equation}

Például, egy $CO$ molekula esetében, az atomok tömege:
\begin{equation}
    ^{12}C: 12 \, u = 1.9926 \times 10^{-26} \, kg \qquad ^{16}O: 16 \, u = 2.6560 \times 10^{-26} \, kg
\end{equation}
És azt is tudjuk a mérések alapján, hogy a $CO$ molekulának a színképvonalak közötti távolság nagyjából $\Delta \nu = 5 \frac{1}{cm}$ vagy Hertz-re átszámolva:
\begin{equation}
    \Delta \nu = 5 \times 10^2 \, m^{-1} \cdot c = 1.5 \times 10^{10} \, Hz
\end{equation}
Ezek alapján a kötési távolság a két atom között:
\begin{equation}
    a = \sqrt{\frac{\hbar}{2 \pi \mu \Delta \nu}} = 2.5 \times 10^{-10} \, m
\end{equation}
Az irodalmi érték ennek a távolságnak $1.128$ angström, ami $1.128 \times 10^{-10} \, m$, tehát a modell elég jó közelítést ad.


\subsubsection{A Hidrogénatom}
A hidrogénatom a legegyszerűbb atom, amely egy protonból és egy elektronból áll. A proton sokkal nehezebb, mint az elektron, így
a proton mozgását elhanyagolhatjuk, és csak az elektron mozgását vizsgáljuk a proton körül (és a protont ezáltal tehetjük az origóba).
A hidrogénatom potenciálja Coulomb-potenciál, ami a következő alakot ölti:
\begin{equation}
    V(r) = - \frac{e^2}{4 \pi \epsilon_0 r}
\end{equation}
Ahol $e$ az elektron töltése, és $\epsilon_0$ a vákuum permittivitása. Ezt a potenciált behelyettesítve a radiális Schrödinger-egyenletbe:
\begin{equation}
   - \frac{\hbar^2}{2m_e} \frac{d^2 u}{dr^2} + \left[ - \frac{e^2}{4 \pi \epsilon_0 r} + \frac{\hbar^2 l (l + 1)}{2m_e r^2} \right] u(r) = E u(r)
\end{equation}
Ahol $m_e$ az elektron tömege, és $u(r) = r R(r)$ és a szögletes zárójelben lévő tagok alkotják az effektív potenciált.
A Coulomb-potenciál maga megengedi a folytonos állapotokat is (E > 0), ami az elektron-proton szóródásokat írja le. Viszont
jelen esetben a kötött állapotokat (E < 0) vizsgáljuk, ahol az elektron a proton körül kering. 
\begin{figure}[H]
    \centering
    \includegraphics[width=0.4\textwidth]{imgs/3-tetel/hidrogen_atom.png}
    \caption{A Hidrogén atom}
\end{figure}
\begin{figure}[H]
    \centering
    \includegraphics[width=0.4\textwidth]{imgs/3-tetel/hidrogen_atom_potencial.png}
    \caption{A Hidrogén atom effektív potenciálja}
\end{figure}

\subsubsection*{A radiális hullámfüggvény}
Először vezessünk be egy új változót:
\begin{equation}
    \kappa \equiv \sqrt{- \frac{2m_e E}{\hbar^2}}
\end{equation}
Ekkor az egyenlet a következő alakot ölti:
\begin{equation}
   \frac{1}{\kappa^2} \frac{d^2 u}{dr^2} = \left[1 - \frac{m_e e^2}{2\pi \epsilon_0 \hbar^2 \kappa} \frac{1}{\kappa r} + \frac{l(l + 1)}{(\kappa r)^2}\right]u
\end{equation}
Adja magát, hogy vezessünk be két további változót:
\begin{equation}
    \varrho \equiv \kappa r \quad \text{és} \quad \varrho_0 \equiv \frac{m_e e^2}{2\pi \epsilon_0 \hbar^2 \kappa}
\end{equation}
Ezekkel a változókkal az egyenlet a következő alakot ölti:
\begin{equation}
   \frac{d^2 u}{d \varrho^2} = \left[1 - \frac{\varrho_0}{\varrho} + \frac{l(l + 1)}{\varrho^2}\right]u
\end{equation}
Ez a differenciálegyenlet elég komplikált, ezért próbáljunk leválasztani róla változókat. Előszö nézzük meg, hogy hogyan viselkedik
a függvény ha $\varrho \to \infty$-be toljuk, tehát a konstans dominál a zárójelben:
\begin{equation}
   \frac{d^2 u}{d \varrho^2} = u
\end{equation}
Ennek a megoldása a következő alakú:
\begin{equation}
    u(\varrho) \approx A e^{-\varrho} + B e^{\varrho}
\end{equation}
De mivel a hullámfüggvénynek normálhatónak kell lennie, ezért a $B$ együtthatót nullának kell vennünk, így:
\begin{equation}
    u(\varrho) \approx A e^{-\varrho}
\end{equation}
Most nézzük meg az egyenletet kis $\varrho$ értékeknél ($\varrho \to 0$). Ilyenkor a $\frac{l(l + 1)}{\varrho^2}$ tag dominál, így az egyenlet közelítőleg:
\begin{equation}
   \frac{d^2 u}{d \varrho^2} = \frac{l(l + 1)}{\varrho^2} u
\end{equation}
Ennek az egyenletnek a megoldása a következő alakú:
\begin{equation}
    u(\varrho) \approx C \varrho^{l + 1} + D \varrho^{-l}
\end{equation}
De mivel a $\varrho^{-l}$ felrobban $\varrho \to 0$-ban, ezért a $D$ együtthatót nullának kell vennünk, így:
\begin{equation}
    u(\varrho) \approx C \varrho^{l + 1}
\end{equation}
Most, hogy megvannak a határfeltételeink, próbáljuk meg felírni a teljes megoldást ezek alapján:
\begin{equation}
    u(\varrho) = \varrho^{l + 1} e^{-\varrho} v(\varrho)
\end{equation}
Ahol $v(\varrho)$ egy ismeretlen függvény, amit meg kell határoznunk, abban reménykedve, hogy ez egyszerűbb lesz mint a $u(\varrho)$. Kiszámolva a deriváltakat:
\begin{equation}
   \frac{du}{d \varrho} = \varrho^{l} e^{-\varrho} \left[(l + 1 - \varrho)v + \varrho \frac{dv}{d \varrho}\right]
\end{equation}
\begin{equation}
   \frac{d^2 u}{d \varrho^2} = \varrho^{l - 1} e^{-\varrho} \left\{\left[-2l - 2 + \varrho + \frac{l(l+1)}{\varrho}\right]v + 2(l + 1 - \varrho) \frac{dv}{d \varrho} + \varrho \frac{d^2 v}{d \varrho^2}\right\}
\end{equation}
Ez egyelőre nem túl bíztató, de helyettesítsük be az eredeti egyenletbe:
\begin{equation}
    \varrho^{l - 1} e^{-\varrho} \left\{\left[-2l - 2 + \varrho + \frac{l(l+1)}{\varrho}\right]v + 2(l + 1 - \varrho) \frac{dv}{d \varrho} + \varrho \frac{d^2 v}{d \varrho^2}\right\} = \left[1 - \frac{\varrho_0}{\varrho} + \frac{l(l + 1)}{\varrho^2}\right] \varrho^{l + 1} e^{-\varrho} v
\end{equation}
\begin{equation}
   \varrho \frac{d^2 v}{d \varrho^2} + 2(l + 1 - \varrho) \frac{dv}{d \varrho} + [\varrho_0 - 2(l + 1)] v = 0
\end{equation}
Végül, tegyük fel, hogy $v(\varrho)$ felírható hatványsor alakban:
\begin{equation}
    v(\varrho) = \sum_{j=0}^{\infty} c_j \varrho^j
\end{equation}
A feladatunk most az, hogy meghatározzuk a $c_j$ együtthatókat. Nézzük meg a deriváltakat:
\begin{equation}
   \frac{dv}{d \varrho} = \sum_{j=0}^{\infty} j c_j \varrho^{j - 1} = \sum_{j=0}^{\infty} (j + 1) c_{j + 1} \varrho^{j}
\end{equation}
Itt a j "dummy indexet" átnevezzük j+1-re, hogy a hatványok megegyezzenek. Ekkor fel lehetne hozni, hogy j=-1-től kéne kezdeni a szummát,
de mivel a (j + 1)=0 tag nullát adna, így ez nem számít. Ugyanezt megtehetjük a második deriválttal is:
\begin{equation}
    \frac{d^2 v}{d \varrho^2} = \sum_{j=0}^{\infty} j (j + 1) c_{j + 1} \varrho^{j - 1}
\end{equation}
visszahelyettesítve:
\begin{equation}
    \begin{aligned}
        & \sum_{j=0}^{\infty} j(j+1)c_{j+1} \varrho^{j} + 2(l + 1) \sum_{j=0}^{\infty} j(j + 1) c_{j + 1} \varrho^{j} - \\
        & - 2 \sum_{j=0}^{\infty} j c_j \varrho^{j} + [\varrho_0 - 2(l + 1)] \sum_{j=0}^{\infty} c_j \varrho^j = 0
    \end{aligned}
\end{equation}
Ez akkor lesz igaz, ha minden együttható nulla:
\begin{equation}
    j(j + 1) c_{j + 1} + 2(l + 1) (j + 1) c_{j + 1} - 2 j c_j + [\varrho_0 - 2(l + 1)] c_j = 0
\end{equation}
Ebből kifejezve az $c_{j + 1}$ együtthatót:
\begin{equation}
    c_{j + 1} = \frac{2(j + l + 1) - \varrho_0}{(j + 1)(j + 2l + 2)} c_j
\end{equation}
Ez egy rekurziós formula, amivel az összes együttható kifejezhető az első együttható segítségével. $c_0$ egy állandó a normalizálás miatt, 
és innen kiindulva az összes többi együttható kiszámolható. Most nézzük meg a sorozat viselkedését nagy j értékeknél. Ilyenkor az együttható közelítőleg:
\begin{equation}
    c_{j + 1} \approx \frac{2j}{j(j+1)}c_j = \frac{2}{j + 1} c_j
\end{equation}
vagyis:
\begin{equation}
    c_j \approx \frac{2^j}{j!} c_0
\end{equation}
Tegyük fel, hogy ez a pontos megoldás, akkor a $v(\varrho)$ függvény a következő alakot ölti:
\begin{equation}
    v(\varrho) = c_0 \sum_{j=0}^{\infty} \frac{(2 \varrho)^j}{j!} = c_0 e^{2 \varrho}
\end{equation}
Ez pedig azt jelenti, hogy a hullámfüggvényünk a következő alakot ölti:
\begin{equation}
    u(\varrho) = \varrho^{l + 1} e^{-\varrho} v(\varrho) = c_0 \varrho^{l + 1} e^{\varrho}
\end{equation}
Tehát a $u(\varrho)$ hullámfüggvény:
\begin{equation}
    u(\varrho) = c_0 \varrho^{l + 1} e^{\varrho}
\end{equation}
Ez felrobban $\varrho \to \infty$-ban, így a hullámfüggvény nem lesz normálható. Ez pontosan az az aszimptotikus viselkedés, amit el szerettünk volna kerülni.
De nem lepődhetünk meg ezen, hiszen ez az eredmény reprezentálja valamilyen aszimptotikus megoldását a radiális egyenletnek, de ezek nem normálhatóak ezért 
nem érdekelnek minket. Mindenesetre ez azt jelenti, hogy a sorozatot le kell vágni egy bizonyos $j$ értéknél, hogy a hullámfüggvény normálható legyen.
Tehát találni kell egy megoldást ahol van egy természetes $N$ szám, amire:
\begin{equation}
    c_{N - 1} \neq 0 \quad \text{de} \quad c_N = 0
\end{equation}
És ezután minden további együttható is nulla lesz. Ebből következik, hogy a rekurziós képletben a számlálónak nullának kell lennie:
\begin{equation}
    2(N + l) - \varrho_0 = 0
\end{equation}
Ahol definiálhatjuk a következő jellemzőt:
\begin{equation}
    n \equiv N + l
\end{equation}
Ekkor a következő összefüggést kapjuk:
\begin{equation}
    \varrho_0 = 2 n \quad \text{ahol } n = 1, 2, 3, \ldots
\end{equation}
Ebből vissza tudjuk számolni az energiákat:
\begin{equation}
    E = - \frac{\hbar^2 \kappa^2}{2m} = - \frac{m_e e^4}{32 \pi^2 \epsilon_0^2 \hbar^2 \varrho_0^2}
\end{equation}
Tehát a megengedett energiák:
\begin{equation}
    E_n = - \frac{m_e e^4}{32 \pi^2 \epsilon_0^2 \hbar^2} \frac{1}{n^2} = \frac{E_1}{n^2} \quad \text{ahol } n = 1, 2, 3, \ldots
\end{equation}
Ez a híres Bohr formula, a kvantummechanika egyik legfontosabb eredménye. Bohr ezt úgy határozta meg, hogy a Schrüdinger-egyenletet még nem ismerték,
és a kvantummechanika még nem létezett mint külön tudományág. Azonban a kvantummechanika megalkotásával ez a formula is igazolást nyert.
Ha pedig megnézzük a $\kappa$ változót:
\begin{equation}
    \kappa = \frac{m_e e^2}{4 \pi \epsilon_0 \hbar^2} \frac{1}{n} = \frac{1}{a_0 n}
\end{equation}
Ahol $a_0$ a Bohr sugár, ami az atomok méretét jellemzi:
\begin{equation}
    a_0 = \frac{4 \pi \epsilon_0 \hbar^2}{m_e e^2} \approx 0.529 \times 10^{-10} \text{ m}
\end{equation}
Innen pedig küvetkezik, hogy:
\begin{equation}
    \varrho = \frac{r}{a_0 n}
\end{equation}
Tehát a térbeli hullámfügvényeket három kvantumszámmal lehet jellemezni (n, l, m):
\begin{equation}
    \psi_{n l m}(r, \theta, \phi) = R_{n l}(r) Y_{l}^m(\theta, \phi)
\end{equation}
Ahol a radiális hullámfüggvény:
\begin{equation}
    R_{n l}(r) = \frac{1}{r}\varrho^{l + 1} e^{-\varrho} v(\varrho)
\end{equation}
Ahol a $v(\varrho)$ egy polinóm $n-l-1$ fokkal $\varrho$-ban, melynek az együtthatói a rekurziós képlettel számolhatók ki:
\begin{equation}
    c_{j + 1} = \frac{2(j + l + 1) - 2n}{(j + 1)(j + 2l + 2)} c_j
\end{equation}
Tehát ha vesszük a legalacsonyabb n értéket (n=1), akkor csak l=0 és m=0 lehetséges → ez az alapállapot.
Az alapállapot hullámfüggvénye tehát:
\begin{equation}
    \psi_{100}(r, \theta, \phi) = R_{10}(r) Y_0^0(\theta, \phi)
\end{equation}
Az energiája pedig:
\begin{equation}
    E_1 = - \frac{m_e e^4}{32 \pi^2 \epsilon_0^2 \hbar^2} \approx -13.6 \text{ eV}
\end{equation}
Ez határozza meg az alapállapotot, ami a "legkisebb" energiát jelöli, amivel az elektront le tudjuk választani a protonról. Ezt kötési energiának is hívjuk.
Ekkor a rekurziós formulában egyszerűsödik az első tag (hisz $j = 0$ miatt $c_1 = 0$ lesz) tehát $v(\varrho)$ egy $c_0$ konstans lesz:
\begin{equation}
    R_ {10}(r) = \frac{1}{r} \varrho^{1} e^{-\varrho} c_0 = c_0 \frac{r}{a_0} e^{-r/a_0} \frac{1}{r} = c_0 \frac{1}{a_0} e^{-r/a_0}
\end{equation}
Ha pedig normalizáljuk ezt, akkor:
\begin{equation}
   \int_{0}^{\infty} |R_{10}(r)|^2 r^2 dr = \frac{|c_0|^2}{a_0^2} \int_{0}^{\infty} e^{-2r/a_0} r^2 dr = |c_0|^2 \frac{a_0}{4} = 1
\end{equation}
Tehát $c_0 = \sqrt{\frac{4}{a_0}}$ és $Y_0^0(\theta, \phi) = \frac{1}{\sqrt{4 \pi}}$. Így az alapállapot teljes hullámfüggvénye:
\begin{equation}
    \psi_{100}(r, \theta, \phi) = \frac{1}{\sqrt{\pi a_0^3}} e^{-r/a_0}
\end{equation}
Ez egy gömbszimmetrikus hullámfüggvény, ami az origó körül koncentrálódik, és exponenciálisan csökken a távolsággal. Ez megfelel annak a fizikai képnek, hogy az elektron a proton körül kering,
és a legvalószínűbb helye az origó közelében van.

Ha $n = 2$ a főkvantumszám, tehát az első gerjesztett állapotban vagyunk, akkor lehet $l = 0$ (ekkor $m = 0$) vagy $l = 1$ (ekkor $m = -1, 0, 1$).
Tehát az $n = 2$ állapotnak összesen 4 féle konfigurációja van, ami az állapot degenerációját jelenti. Az energiája pedig:
\begin{equation}
    E_2 = \frac{E_1}{2^2} = \frac{E_1}{4} \approx -3.4 \text{ eV}
\end{equation}
Ekkor az együtthatók a rekurziós képlet alapján kiszámolhatók, és a radiális hullámfüggvények is meghatározhatók. Például ha $l = 0$ és $m = 0$, akkor:
\begin{equation}
    c_1 = \frac{2(0 + 0 + 1) - 4}{(0 + 1)(0 + 0 + 2)} c_0 = - \frac{2}{2} c_0 = - c_0
\end{equation}
\begin{equation}
    c_2 = 0 \quad \text{(mivel a sorozatot le kell vágni $N= n - l = 1$-nél)}
\end{equation}
Szóval $v(\varrho) = c_0 (1 - \varrho)$, tehát:
\begin{equation}
    R_{20}(r) = \frac{c_0}{2a} (1 - \frac{r}{2a_0}) e^{-r/2a_0}
\end{equation}
Ha pedig $l = 1$ és $m = 0$, akkor:
\begin{equation}
    c_1 = \frac{2(0 + 1 + 1) - 4}{(0 + 1)(0 + 2 + 2)} c_0 = \frac{0}{4} c_0 = 0
\end{equation}
\begin{equation}
    c_2 = 0 \quad \text{(mivel a sorozatot le kell vágni $N= n - l = 1$-nél)}
\end{equation}
Szóval $v(\varrho) = c_0$, tehát:
\begin{equation}
    R_{21}(r) = \frac{c_0}{4 a_0^2} r e^{-r/2a_0}
\end{equation}
A $c_0$ minden esetben normalizálással határozható meg. Láthatjuk tehát, hogy egy adott $n$ értékre $l = 0, 1, 2, \dots, n-1$ lehetséges, és minden $l$ értékre $m = -l, -l+1, \dots, l-1, l$ lehetséges.
Tehát egy adott energiaszint degenerációja:
\begin{equation}
    d(n) = \sum_{l=0}^{n-1} (2l + 1) = n^2
\end{equation}
\begin{figure}[H]
    \centering
    \includegraphics[width=0.6\textwidth]{imgs/3-tetel/hidrogen_atom_energiak.png}
    \caption{A Hidrogén atom energiaszintjei és hullámfüggvényei; $n = 1$ az alapállapot $E_1 = -13.6 \text{ eV}$ enegiával;
    Végtelen számú állapot fér be $n = 5$ és $n = \infty$ közé; $E_\infty = 0$ választja el a kötött és szabad állapotokat.}
\end{figure}

A $v(\varrho)$ polinóm amit használtunk a megoldásban a matematikusok számára már jól ismert asszociált Laguerre polinómok normalizálás nélkül:
\begin{equation}
    v(\varrho) = L_{n - l - 1}^{2l + 1}(2\varrho)
\end{equation}
Ahol:
\begin{equation}
    L_p^q(x) = (-1)^q \left(\frac{d}{dx}\right)^q L_{p + q}(x)
\end{equation}
az Asszociált Laguerre polinóm és
\begin{equation}
    L_q(x) = \frac{e^x}{q!} \left(\frac{d}{dx}\right)^q (x^q e^{-x})
\end{equation}
A $q$-edik Laguerre polinóm. Innen a normalizált radiális hullámfüggvénye a Hidrogén atomnak:
\begin{equation}
    \psi_{n l m} = \sqrt{\left(\frac{2}{n a_0}\right)^3 \frac{(n - l - 1)!}{2 n [(n + l)!]^3}} e^{-r/na_0} \left(\frac{2r}{n a_0}\right)^l L_{n - l - 1}^{2l + 1} \left(\frac{2r}{n a_0}\right) Y_l^m(\theta, \phi)
\end{equation}
Ami nem egy túl szép megoldás, de a kevés, zárt alakban megoldható valós rendszer megoldásának egyike. A hullámfüggvények kölcsönösen ortogonálisak és normálhatók:
\begin{equation}
    \int \psi_{n' l' m'}^* (r, \theta, \phi) \psi_{n l m} (r, \theta, \phi) r^2 \sin(\theta) dr d\theta d\phi = \delta_{n n'} \delta_{l l'} \delta_{m m'}
\end{equation}
Ez abból következik, hogy a gömbfüggvények ortogonálisak és ($n \neq n'$) és abból, hogy ezek a Hamilton-operátor sajátfüggvényei különböző sajátértékekhez.

\subsubsection*{Hidrogén atom spektruma}
Elméletileg, ha a hidrogén atomot egy $\psi_{n l m}$ állapotba rakjuk, akkor ott is marad örökké. Ha viszont egy kicsit "megpiszkáljuk" az atomot (pl. egy fotonnal ütköztetjük),
akkor az atom gerjesztett állapotba kerülhet ($n > 1$) vagy visszaeshet egy alacsonyabb állapotba. Ezek a perturbációk a valóságban mindig jelen vannak, így az atom sosem marad egy állapotban,
időnként feljegb ugrik egyet majd visszaesik (ezeket kvantum ugrásoknak is nevezik). Amikor visszaesik egy alacsonyabb állapotba, akkor meg kell szabadulnia a fölösleges energiától,
amit általában elektromágneses sugárzás formájában tesz meg (foton kibocsátás). Ennek a kibocsájtott fotonnak pontosan annyi lesz az energiája, mint a két állapot közötti energiakülönbség:
\begin{equation}
    E_{\text{foton}} = E_i - E_f = -13.6 \text{ eV} \left(\frac{1}{n_i^2} - \frac{1}{n_f^2}\right)
\end{equation}
Azt pedig a Planck összefüggésből tudjuk, hogy a foton energiája és frekvenciája között a következő összefüggés áll fenn:
\begin{equation}
    E_{\text{foton}} = h \nu
\end{equation}
És a hullámhossz és a frekvencia között pedig:
\begin{equation}
    c = \lambda \nu
\end{equation}
Ebből következik, hogy a kibocsájtott foton hullámhossza:
\begin{equation}
    \frac{1}{\lambda} = \mathcal{R} \left(\frac{1}{n_f^2} - \frac{1}{n_i^2}\right)
\end{equation}
Ahol $\mathcal{R}$ a Rydberg állandó, ami a következő alakot ölti:
\begin{equation}
    \mathcal{R} \equiv \frac{m_e}{4 \pi c \hbar^3} \left(\frac{e^2}{4 \pi \epsilon}\right)^2 \approx 1.097 \times 10^7 \text{ m}^{-1}
\end{equation}
Ez pedig a híres Rydberg formula, amit empirikus úton fedeztek fel a hidrogén spektrumának vizsgálatakor. Itt $n_i$ az induló állapot főkvantumszáma,
és $n_f$ a végállapot főkvantumszáma. Különböző $n_f$ értékekhez különböző sorozatok tartoznak:
\begin{itemize}
    \item Lyman sorozat: $n_f = 1$, ultraibolya tartomány
    \item Balmer sorozat: $n_f = 2$, látható tartomány
    \item Paschen sorozat: $n_f = 3$, infravörös tartomány
    \item Brackett sorozat: $n_f = 4$, infravörös tartomány
    \item Pfund sorozat: $n_f = 5$, infravörös tartomány
    \item Humphreys sorozat: $n_f = 6$, infravörös tartomány
\end{itemize}
\begin{figure}[H]
    \centering
    \includegraphics[width=0.6\textwidth]{imgs/3-tetel/hidrogen_atom_spektrum.png}
    \caption{A Hidrogén atom spektrumának Energia szintjei és átmenetei.}
\end{figure}


\end{document}